% ============================================================
%  第三章：从偏好到效用（精讲版）
%  编译：LuaLaTeX（两遍）
% ============================================================
% 第 3 章 统一书稿前言包装
\newcommand{\lectureno}{第三章}
\newcommand{\lectureshort}{从偏好到效用}
% ============================================================
%  微观经济学理论课程讲义 —— 统一课程前言（四讲共用）
%  版式参考 Nolan Miller, Notes on Microeconomic Theory (SWP/notes3 风格)
%  编译：LuaLaTeX（两遍）
%  各讲 preamble.tex 只需设定 \lectureno / \lectureshort 后 \input 本文件。
% ============================================================
\documentclass[11pt,a4paper,fontset=none,oneside,openany]{ctexbook}

% ---------- 版式 ----------
\usepackage{geometry}
\geometry{
  a4paper,
  top=2.8cm, bottom=2.8cm,
  left=2.6cm, right=2.6cm,
  headsep=0.85cm, footskip=1.25cm
}
\setlength{\headheight}{14pt}

% ---------- 字体（全部矢量嵌入，保证 PDF 高清晰度）----------
% 注意：本机 macOS 已移除 Kaiti SC，CJK 斜体统一回退到 Songti SC；
%        Songti SC / Heiti SC 经 OSFONTDIR 指向 Supplemental 后由 luaotfload 解析。
\usepackage{fontspec}
\setmainfont{TeX Gyre Termes}[Ligatures=TeX]
\setsansfont{TeX Gyre Heros}[Ligatures=TeX,Scale=MatchLowercase]
\setmonofont{TeX Gyre Cursor}[Scale=MatchLowercase]
\setCJKmainfont{Songti SC}[BoldFont={Songti SC Bold},ItalicFont={Songti SC}]
\setCJKsansfont{Heiti SC}[BoldFont={Heiti SC}]

% ---------- 数学 ----------
\usepackage{amsmath}
\usepackage{mathtools}
\usepackage{amsthm}
\usepackage{unicode-math}
\setmathfont{TeX Gyre Termes Math}
\allowdisplaybreaks[2]
\setlength{\abovedisplayskip}{9pt plus 3pt minus 3pt}
\setlength{\belowdisplayskip}{9pt plus 3pt minus 3pt}

% ---------- 定理类环境（书稿格式：按章编号，如“定理 1.1”）----------
\theoremstyle{plain}
\newtheorem{theorem}{定理}[chapter]
\newtheorem{proposition}[theorem]{命题}
\newtheorem{lemma}[theorem]{引理}
\newtheorem{corollary}[theorem]{推论}
\theoremstyle{definition}
\newtheorem{definition}[theorem]{定义}
\newtheorem{condition}[theorem]{条件}
\newtheorem{example}[theorem]{例}
\theoremstyle{remark}
\newtheorem{remark}[theorem]{注}
% 公式、图、表按章编号（书稿格式：式 (1.1)、图 1.1、表 1.1）
\numberwithin{equation}{chapter}
\numberwithin{figure}{chapter}
\numberwithin{table}{chapter}

% 定理类环境标题后的标点：amsthm 默认加英文 "."，此处统一改为中文句号。
% 原理：amsthm 的 \@thm 每次开环境都执行 \thm@headpunct{.}，而样式宏在其之后展开，
% 故把 \thm@headpunct{。} 追加到三个样式宏末尾即可覆盖（用 \AtBeginDocument 以免被 ctex 覆盖）。
\makeatletter
\AtBeginDocument{%
  \let\WBth@plain\th@plain
  \let\WBth@definition\th@definition
  \let\WBth@remark\th@remark
  \def\th@plain{\WBth@plain\thm@headpunct{。}}%
  \def\th@definition{\WBth@definition\thm@headpunct{。}}%
  \def\th@remark{\WBth@remark\thm@headpunct{。}}%
}
\makeatother

% 无编号小标题（仿 notes3 的 “A Note on ...”）
\newcommand{\notehead}[1]{\par\medskip\noindent{\bfseries #1}\par\smallskip}

% ---------- 表格 ----------
\usepackage{booktabs}
\usepackage{array}
\renewcommand{\arraystretch}{1.25}
\setlength{\tabcolsep}{8pt}

% ---------- 图形 ----------
\usepackage{graphicx}
\usepackage{tikz}
\usepackage{pgfplots}
\pgfplotsset{compat=1.18}
\usetikzlibrary{arrows.meta,positioning,shapes.geometric,calc,patterns,decorations.pathreplacing,decorations.markings}
\usepackage{float}
\usepackage{caption}
\captionsetup{font=small,labelfont=bf,labelsep=quad,skip=4pt}

% 统一图形配色（黑白模式：黑/灰梯度，供所有 TikZ 图复用）
% —— 全局黑白化：原彩色调色板统一改为黑/灰，保证打印与投影均可读 ——
\definecolor{AxisBlue}{RGB}{0,0,0}        % 轴线 / 主曲线：黑
\definecolor{SoftBlue}{RGB}{130,130,130}  % 节点填充：中灰
\definecolor{LightRed}{RGB}{224,224,224}  % 浅填充：浅灰
\definecolor{DeepBlue}{RGB}{0,0,0}        % 描边：黑
\definecolor{MidBlue}{RGB}{70,70,70}      % 次曲线 / 网格：深灰
\tikzset{
  >=Stealth,
  every picture/.append style={font=\small},
  axisline/.style={thick,color=black,-Stealth},
  axislabel/.style={color=black},
  fillsoft/.style={fill=LightRed,fill opacity=0.6,draw=black,thick},
  nodeblue/.style={circle,fill=black!85,draw=black,thick,inner sep=2pt},
  nodewhite/.style={circle,fill=white,draw=black,thick,inner sep=2pt},
  arrblue/.style={->, thick, color=black},
  boxsoft/.style={draw=black, rounded corners=2pt, align=center, inner sep=5pt, fill=black!6},
  % 多曲线 / 多集合区分：以线型 + 标记代替颜色
  curveA/.style={thick,black,solid},
  curveB/.style={thick,black,dashed},
  curveC/.style={thick,black,dotted},
  markA/.style={mark=*,mark size=2.2pt,color=black},
  markB/.style={mark=square*,mark size=2.0pt,color=black},
  markC/.style={mark=triangle*,mark size=2.4pt,color=black},
  % 图内区域 / 点标签（白底避免与图形重叠）
  figlabel/.style={font=\scriptsize,fill=white,inner sep=1pt,outer sep=1pt,text=black},
}
% 图形说明（黑白模式下对图形区域 / 曲线 / 关键点的解释，置于图形下方）
\newcommand{\figexplan}[1]{%
  \par\smallskip\noindent{\small\itshape 图形说明：#1}\par\smallskip}

% ---------- 页眉页脚（书稿格式：左=书名，右=章次与章名）----------
\usepackage{fancyhdr}
% 章标题：居中、大号（书稿章首样式）
\CTEXsetup[format={\centering\LARGE\bfseries}]{chapter}
\CTEXsetup[beforeskip=16pt,afterskip=14pt]{chapter}
\CTEXsetup[format={\large\bfseries}]{section}
\CTEXsetup[format={\normalsize\bfseries}]{subsection}
\pagestyle{fancy}
\fancyhf{}
\providecommand{\booktitle}{微观经济学理论：可行集、偏好与效用}
\providecommand{\lectureno}{第一章}
\providecommand{\lectureshort}{课程讲义}
\fancyhead[L]{\small\itshape \booktitle}
\fancyhead[R]{\small \lectureno\quad \lectureshort}
\fancyfoot[C]{\small \thepage}
\renewcommand{\headrulewidth}{0.4pt}
\fancypagestyle{plain}{\fancyhf{}\fancyfoot[C]{\small \thepage}\renewcommand{\headrulewidth}{0pt}}

% ---------- 正文细节 ----------
\usepackage{setspace}
\setstretch{1.12}
\usepackage{microtype}
\tolerance=900
\emergencystretch=2em
\raggedbottom

% ---------- 超链接与 PDF 书签 ----------
\usepackage{hyperref}
\hypersetup{
  colorlinks=true,
  linkcolor=black,
  urlcolor=black,
  citecolor=black,
  pdfencoding=auto,
  bookmarks=true,
  pdftitle={\booktitle\quad \lectureno：\lectureshort},
  pdfsubject={微观经济学理论书稿}
}
% 章节标题中常见数学符号的安全书签映射（避免 "Improper alphabetic constant"）
\pdfstringdefDisableCommands{%
  \def\leftrightarrow{<->}%
  \def\star{*}%
  \def\S{S}%
  \def\succeq{\text{≽}}%
  \def\succ{\text{≻}}%
  \def\sim{\text{∼}}%
}

% ---------- 常用记号 ----------
\newcommand{\R}{\mathbb{R}}
\newcommand{\N}{\mathbb{N}}
\newcommand{\Z}{\mathbb{Z}}
\newcommand{\Q}{\mathbb{Q}}
\newcommand{\succeqq}{\succeq}
\newcommand{\argmax}{\operatorname*{arg\,max}}
\newcommand{\argmin}{\operatorname*{arg\,min}}
\newcommand{\st}{\text{s.t.}\ }

\setcounter{chapter}{2}


\begin{document}

\chapter{从偏好到效用：效用表示的存在性、构造与性质对应}

\noindent\textbf{本章任务。}第二章建立了偏好公理（完备性、传递性）与显示偏好理论，
遗留一个核心问题：\textbf{什么样的偏好可以被一个效用函数“数字化”？}
本章给出完整回答，沿四个层次推进：
\begin{enumerate}
  \item \textbf{存在性}：可表示偏好的必要条件（理性 + 可数序稠密子集），
        词典序偏好作为“理性但不可表示”的基准反例（§3.1）；
  \item \textbf{构造}：有限集上的归纳构造与排名构造（§3.2），
        无限集上的权数构造、连续编号法与缺口分析（§3.3--§3.4）；
  \item \textbf{性质对应}：偏好的结构特性（单调、凸、位似、可分、拟线性、连续）
        如何一一翻译为效用函数的解析性质（§3.5），
        可分性与拟线性的完整表示定理及其证明（§3.6--§3.7）；
  \item \textbf{边界}：实验证据对公理体系的系统性偏离，
        每一类行为效应精确对应哪一条公理被破坏（§3.8）。
\end{enumerate}

\noindent 全章的枢纽是三组\textbf{充要条件}：
\begin{gather*}
\boxed{\ \text{可表示}\ \Longleftrightarrow\ \text{理性}+\text{可数序稠密子集}\ }\\[2pt]
\boxed{\ \text{可加表示}\ \Longleftrightarrow\ \text{坐标独立}+\text{矩形条件}\ }\\[2pt]
\boxed{\ \text{拟线性表示}\ \Longleftrightarrow\ \text{计价物三条件}\ }
\end{gather*}

\medskip
\noindent\textbf{阅读提示}：§3.4 的连续编号构造是全章技术核心，
§3.6、§3.7 的两个表示定理是后续消费者理论、福利分析的工具库入口；
带 $\star$ 的小节为加深内容，第一遍可略读。

\tableofcontents
\clearpage

% ============================================================
\section{哪些偏好可以被效用表示}
% ============================================================

\subsection{效用表示的定义与序数等价}

\begin{definition}[效用表示（回顾 \S2.2）]\label{def:rep}
设 $\succeq$ 为备选集 $X$ 上的弱偏好关系。函数 $u:X\to\R$ \textbf{表示}（represents）$\succeq$，
当且仅当
\[
\forall x,y\in X:\qquad x\succeq y\ \Longleftrightarrow\ u(x)\ge u(y).
\]
此时称 $\succeq$ \textbf{可效用表示}（utility-representable），$u$ 称为 $\succeq$ 的一个效用函数。
\end{definition}

表示定义是\textbf{双向}的：不仅“偏好高 $\Rightarrow$ 数值高”，而且“数值高 $\Rightarrow$ 偏好高”。
因此 $u(x)=u(y)\iff x\sim y$，$u(x)>u(y)\iff x\succ y$：
一个效用函数同时编码了 $\succeq$、$\succ$、$\sim$ 三层信息。

\begin{proposition}[可表示 $\Rightarrow$ 理性]\label{prop:rep-rational}
若 $\succeq$ 可被某 $u:X\to\R$ 表示，则 $\succeq$ 完备且传递（即为理性偏好/弱序）。
\end{proposition}

\begin{proof}
\textbf{完备性}：任取 $x,y\in X$，$u(x),u(y)\in\R$。实数的全序性给出 $u(x)\ge u(y)$ 或 $u(y)\ge u(x)$，
由表示定义即 $x\succeq y$ 或 $y\succeq x$。

\textbf{传递性}：设 $x\succeq y$、$y\succeq z$。由表示定义 $u(x)\ge u(y)\ge u(z)$，
实数 $\ge$ 传递给出 $u(x)\ge u(z)$，即 $x\succeq z$。
\end{proof}

命题~\ref{prop:rep-rational} 给出第一道门槛：\textbf{非理性偏好（不完备或不传递）一律不可表示}。
但理性远不是充分条件——§\ref{subsec:lex} 将给出理性却不可表示的偏好。

\begin{proposition}[序数等价：表示不唯一]\label{prop:ordinal}
设 $u$ 表示 $\succeq$，$f:\R\to\R$ 严格递增，则 $\tilde u\coloneqq f\circ u$ 也表示 $\succeq$；
反之，若 $u,\tilde u$ 表示同一 $\succeq$，则存在严格递增 $f:u(X)\to\R$ 使 $\tilde u=f\circ u$。
\end{proposition}

\begin{proof}
\textbf{正向}：$x\succeq y\iff u(x)\ge u(y)\iff f(u(x))\ge f(u(y))\iff \tilde u(x)\ge \tilde u(y)$，
第二步用 $f$ 严格递增（保持 $\ge$ 的双向性）。

\textbf{反向}：在值域 $u(X)$ 上定义 $f(t)=\tilde u\bigl(u^{-1}(t)\bigr)$（任取原像）：
若 $u(x)=u(x')$，则 $x\sim x'$，故 $\tilde u(x)=\tilde u(x')$，$f$ 良定义；
若 $t_1>t_2$，取 $u(x_1)=t_1>u(x_2)=t_2$，则 $x_1\succ x_2$，故 $\tilde u(x_1)>\tilde u(x_2)$，即 $f$ 严格递增。
\end{proof}

\begin{remark}[序数性质与基数性质]\label{rem:cardinal}
命题~\ref{prop:ordinal} 意味着效用数值本身\textbf{没有绝对意义}：$u(x)=10$ 与 $u(x)=2$ 之差为 $8$
不代表“好 8 个单位”，因为 $\tilde u=\sqrt{u}$ 同样是合法表示。
凡在\textbf{一切严格递增变换下保持不变}的性质称为\textbf{序数性质}
（如“无差异曲线的形状”“最优选择集”）；
只对特定表示成立的性质称为\textbf{基数性质}（如“边际效用递减”）。
§\ref{subsec:ordinal-cardinal} 将严格区分这两类性质——这是性质对应问题（§3.5）的逻辑前提。
\end{remark}

\begin{figure}[H]
\centering
\begin{tikzpicture}[font=\scriptsize,scale=1.18,transform shape]
  % ---------- (a) 定义域离散：序上存在相邻对 ----------
  \draw[thick,fill=gray!15] (0,0) rectangle (3.6,3.0);
  \node[font=\small\bfseries] at (1.8,2.62) {(a) 离散定义域：序上存在相邻对};
  \draw[thick,dashed] (0.5,1.45) -- (3.1,1.45);
  \node[figlabel,anchor=east] at (0.42,1.45) {$X$};
  \fill (0.9,1.45) circle (1.8pt) node[below=3pt] {$x$};
  \fill (2.7,1.45) circle (1.8pt) node[below=3pt] {$y$};
  \draw[decorate,decoration={brace,amplitude=4pt,mirror}] (1.05,2.00) -- (2.55,2.00)
    node[midway,above=3pt,font=\scriptsize] {此处无任何元素};
  \node[align=center,text width=3.3cm,font=\scriptsize] at (1.8,0.62)
    {虚线只表示\emph{序的位置}，不代表区间内的点。\\ $x$ 与 $y$ 是相邻的无差异类，\\
     $u(x)$ 与 $u(y)$ 之间留下一段实数空隙};
  % ---------- (b) 序稠密：任意一对之间都有元素 ----------
  \begin{scope}[xshift=6.2cm]
    \draw[thick,fill=gray!15] (0,0) rectangle (3.6,3.0);
    \node[font=\small\bfseries] at (1.8,2.62) {(b) 序稠密：任意一对之间有元素};
    \draw[thick,dashed] (0.5,1.45) -- (3.1,1.45);
    \node[figlabel,anchor=east] at (0.42,1.45) {$X$};
    \fill (0.7,1.45) circle (1.8pt) node[below=3pt] {$x$};
    \fill (1.5,1.45) circle (1.8pt) node[above=3pt] {$d$};
    \fill (2.3,1.45) circle (1.8pt) node[above=3pt] {$d'$};
    \fill (3.1,1.45) circle (1.8pt) node[below=3pt] {$y$};
    \draw[decorate,decoration={brace,amplitude=4pt}] (0.9,1.90) -- (2.9,1.90)
      node[midway,above=3pt,font=\scriptsize] {每个空隙都插入了 $d,d'$};
    \node[align=center,text width=3.3cm,font=\scriptsize] at (1.8,0.62)
      {对每一对 $x\succ y$ 都存在 $d$ 夹在中间，\\ 于是可以把 $u(d)$ 插进 $u(x)$ 与 $u(y)$ 之间，\\
       序被填满、跳跃被细分};
  \end{scope}
\end{tikzpicture}
\caption{可数序稠密子集为何必要：序上的缺口与稠密}
\end{figure}
\figexplan{两幅图的横轴都是\emph{序的位置}（虚线），不是实数轴本身——这一点是读图的关键。左面板的定义域是\textbf{离散}的：$x$ 与 $y$ 是相邻的两个无差异类，它们之间没有任何方案，因此无论怎样定义 $u$，$u(x)$ 与 $u(y)$ 之间总会留下一段实数空隙，而空隙中的值没有任何方案与之对应，序无法被完整编码。右面板表示序稠密：每一对相邻方案之间都插入了 $d,d'$，空隙被不断细分，跳跃至多可数个，从而序可以被 $\R$ 上的数值完整承载。}

\begin{theorem}[可表示的必要条件]\label{thm:dense-necessary}
若弱序 $\succeq$ 可被效用函数表示，则 $X$ 中存在\textbf{至多可数}的序稠密子集。
\end{theorem}

\begin{proof}
设 $u$ 表示 $\succeq$。分两类收集“见证点”构成 $Z$。

\textbf{第一类：跳跃之间的点。}考虑 $u(X)\subseteq\R$ 中所有满足如下条件的无差异类 $[x]$：
存在无差异类 $[y]$ 使 $u(x)>u(y)$ 且开区间 $\bigl(u(y),u(x)\bigr)$ 内\textbf{没有} $u$ 的值
（即 $[x]$ 与 $[y]$ 之间是一个“跳跃”）。
对每个这样的跳跃，区间 $\bigl(u(y),u(x)\bigr)$ 是 $\R$ 中非空开区间，
且不同跳跃对应的区间\textbf{两两不交}（每个跳跃区间以值域中的两个相邻值为端点）。
实数轴上任何一族两两不交的非空开区间至多可数
（每个区间任取一有理数，即得该族到 $\Q$ 的单射），
故跳跃点至多可数个；从每个跳跃两端的无差异类中各取一代表元，记入 $Z_1$。

\textbf{第二类：非跳跃点的逼近。}对不是跳跃端点的无差异类 $[x]$，
若 $[x]$ 上方（或下方）紧邻没有跳跃，则 $u(x)$ 可被值域中的其他值从一侧逼近。
对每个 $n\in\N$，在 $u(X)$ 的每个长度落入 $\bigl[\frac{1}{n+1},\frac{1}{n}\bigr)$ 的
“空洞”端点无差异类中取代表元；同样的不交区间论证给出至多可数个，记为 $Z_2$。

令 $Z=Z_1\cup Z_2$，则 $Z$ 至多可数。任取 $x\succ y$：
若 $u(x)$ 与 $u(y)$ 之间隔着某个跳跃，跳跃端点的代表元即满足序稠密要求；
否则 $u(X)$ 在 $(u(y),u(x))$ 内有值（无跳跃），任取 $u(z)\in\bigl(u(y),u(x)\bigr)$ 对应的无差异类，
其代表元可由 $Z_2$ 中的逼近点替代满足要求。故 $Z$ 序稠密。
\end{proof}

定理~\ref{thm:dense-necessary} 的逆命题也成立（充分性），其证明就是 §3.4 连续编号构造的目标：
\[
\boxed{\ \succeq\ \text{可表示}\ \Longleftrightarrow\ \succeq\ \text{为弱序，且}\ X\ \text{含可数序稠密子集}\ }
\]
这是可表示性问题的\textbf{完全刻画}：没有任何额外的拓扑假设，纯粹的序条件。

\subsection{基准反例：词典序偏好不可表示}\label{subsec:lex}

\begin{definition}[词典序偏好]\label{def:lex}
在 $X=\R^2_+$ 上定义 $\succeq_L$：
\[
(x_1,x_2)\succeq_L (y_1,y_2)\ \Longleftrightarrow\
\Bigl[x_1>y_1\ \ \text{或}\ \ \bigl(x_1=y_1\ \text{且}\ x_2\ge y_2\bigr)\Bigr].
\]
即先比较第一坐标；第一坐标相同时才比较第二坐标——如同字典先按首字母排序。
\end{definition}

\begin{proposition}[$\succeq_L$ 是理性偏好]\label{prop:lex-rational}
$\succeq_L$ 完备、传递（从而是弱序），且满足强单调性：$x\gg y\Rightarrow x\succ_L y$。
\end{proposition}

\begin{proof}
\textbf{完备性}：任取 $x,y$。若 $x_1\ne y_1$，第一坐标大者严格优；若 $x_1=y_1$，比较 $x_2,y_2$ 即得结论。
\textbf{传递性}：设 $x\succeq_L y$、$y\succeq_L z$。
若 $x_1>y_1$ 或 $y_1>z_1$，则 $x_1>z_1$，故 $x\succ_L z$；
若 $x_1=y_1=z_1$，则 $x_2\ge y_2\ge z_2$，故 $x\succeq_L z$。
\textbf{强单调性}：$x\gg y$ 蕴含 $x_1>y_1$ 或（$x_1\ge y_1$ 且 $x_2>y_2$ 当 $x_1=y_1$），两种情形均得 $x\succ_L y$。
\end{proof}

\begin{center}
\begin{tikzpicture}[scale=1.24,transform shape]
  % axes
  \draw[axisline] (0,0) -- (6.4,0) node[below left, axislabel] {$x_1$};
  \draw[axisline] (0,0) -- (0,4.4) node[below left, axislabel] {$x_2$};
  % vertical first-coordinate iso-lines
  \foreach \x in {1,2,3,4,5} {\draw[MidBlue!35, thin] (\x,0) -- (\x,4.1);}
  % point and better/worse sets
  \node[nodeblue, label=below right:{$x$}] at (3,2.2) {};
  \draw[arrblue] (3,2.2) -- (3,3.9) node[midway, right] {$y\succ_L x$（同一竖线向上）};
  \draw[arrblue] (3,2.2) -- (3,0.25) node[midway, right] {$y\prec_L x$（同一竖线向下）};
  \draw[arrblue] (3.05,2.2) -- (5.7,2.2) node[pos=0.82, below=3pt,fill=white,inner sep=1pt] {右侧点严格优于 $x$};
  \draw[arrblue] (2.95,2.1) -- (0.35,2.1) node[pos=0.62, above=3pt,fill=white,inner sep=1pt] {左侧点严格劣于 $x$};
  \node[align=center,text width=6.4cm,fill=white,inner sep=2pt] at (3,-0.9) {无差异类 $=\{x\}$：无无差异曲线；\\只有竖直方向和左右方向的严格比较};
\end{tikzpicture}
\end{center}
\figexplan{横轴 $x_1$、纵轴 $x_2$。每根浅灰竖线是“第一坐标固定”的比较基准：在同一条竖线上，位置越高越优、越低越劣；整条竖线向右移动严格变优、向左移动严格变劣。由于词典序先比较 $x_1$，无差异类退化为单点 $\{x\}$，故不存在通常意义的无差异曲线，只有竖直方向与左右方向两组严格比较。}

\begin{theorem}[词典序偏好不可效用表示]\label{thm:lex}
不存在任何 $u:\R^2_+\to\R$ 表示 $\succeq_L$。
\end{theorem}

\begin{proof}
反设 $u$ 表示 $\succeq_L$。对每个 $t\in\R_+$，比较两点 $(t+1,0)$ 与 $(t,1)$：
\[
(t+1,0)\succ_L(t,1)\quad(\text{第一坐标 }t+1>t)
\ \Longrightarrow\ u(t+1,0)>u(t,1).
\]
于是对每个 $t\in\R_+$，开区间
\[
I_t\ \coloneqq\ \bigl(u(t,0),\,u(t,1)\bigr)
\]
非空（因 $(t,1)\succ_L(t,0)$，第二坐标 $1>0$，故 $u(t,1)>u(t,0)$）。

\textbf{关键一步：$\{I_t\}_{t\in\R_+}$ 两两不交。}
取 $s<t$。则 $(s,1)\prec_L(t,0)$（第一坐标 $s<t$），故 $u(s,1)<u(t,0)$，即
\[
I_s=\bigl(u(s,0),u(s,1)\bigr)\ \text{整体位于}\ u(t,0)\ \text{左侧},\qquad
I_s\cap I_t=\varnothing.
\]

每个 $I_t$ 是非空开区间，任取有理数 $q(t)\in I_t\cap\Q$。
由两两不交，$t\ne s\Rightarrow q(t)\ne q(s)$，故 $q:\R_+\to\Q$ 是\textbf{单射}。
但这要求 $|\R_+|\le|\Q|$，与 $\R_+$ 不可数、$\Q$ 可数矛盾。故表示不存在。
\end{proof}

\begin{remark}[反例的结构]\label{rem:lex-structure}
证明的三步骨架值得记住：(i) 对每个“层级”$t$ 制造一个非空开区间 $I_t$；
(ii) 用偏好的序结构证明这些区间两两不交；(iii) 用“不交开区间族至多可数”完成基数矛盾。
词典序偏好有\textbf{不可数多个}无差异类（每个单点都是一类），且类与类之间都隔着“跳跃”，
直接违反定理~\ref{thm:dense-necessary} 的必要条件——可数序稠密子集不存在：
任何序稠密子集必须在每个层级 $t$ 内部布点，故至少与 $\R_+$ 等势。
\end{remark}

\subsection{本章路线图}

\begin{center}
\begin{tikzpicture}[
  font=\small,
  box/.style={draw=AxisBlue, rounded corners=2pt, align=center, inner sep=5pt, fill=MidBlue!10},
  arr/.style={-{Stealth[length=2.5mm]}, thick, color=AxisBlue},
  node distance=0.9cm and 1.15cm
]
\node[box] (ax) {偏好公理\\ 完备 + 传递};
\node[box, right=of ax] (fin) {有限集\\ 归纳 / 排名构造\\ （§3.2）};
\node[box, right=of fin] (inf) {无限集\\ 权数构造（可数）\\ 连续编号（连续统）\\ （§3.3--§3.4）};
\node[box, below=of ax] (nsec) {必要条件\\ 可数序稠密子集\\ 词典序反例（§3.1）};
\node[box, below=of fin] (prop) {性质对应\\ 单调/凸/位似/可分\\ 拟线性/连续（§3.5）};
\node[box, below=of inf] (rep2) {结构表示定理\\ 可加表示（§3.6）\\ 拟线性表示（§3.7）};
\node[box, below=of prop] (beh) {行为批判\\ 公理的系统性失效（§3.8）};
\draw[arr] (ax) -- (fin);
\draw[arr] (fin) -- (inf);
\draw[arr] (ax) -- (nsec);
\draw[arr] (nsec) -- (prop);
\draw[arr] (inf) -- (rep2);
\draw[arr] (prop) -- (rep2);
\draw[arr] (prop) -- (beh);
\draw[arr] (rep2) -- (beh);
\end{tikzpicture}
\end{center}
\figexplan{上支为构造主线：偏好公理 $\to$ 有限集构造 $\to$ 无限集构造；下支为约束与性质：必要条件（可数序稠密子集，词典序反例）与性质对应 $\to$ 结构表示定理（可加表示、拟线性表示）；两支最终汇入行为批判。箭头表示逻辑推进方向，方框内注明对应小节。}

% ============================================================
\section{有限集上效用函数的存在性}
% ============================================================

\begin{figure}[H]
\centering
\begin{tikzpicture}[font=\scriptsize,scale=1.12,transform shape]
  \draw[thick,fill=gray!15] (0,0) rectangle (1.9,2.6);
  \node[font=\small\bfseries] at (0.95,2.28) {$n=1$};
  \fill (0.95,1.45) circle (1.5pt) node[right=3pt] {$a$};
  \node[align=center] at (0.95,0.78) {$u(a)=1$};
  \node[align=center] at (0.95,0.25) {起点：单元素};
  \draw[->,thick] (1.98,1.3) -- (2.42,1.3);
  \draw[thick,fill=gray!15] (2.5,0) rectangle (4.4,2.6);
  \node[font=\small\bfseries] at (3.45,2.28) {$n=2$};
  \fill (3.10,1.45) circle (1.5pt) node[left=3pt] {$a$};
  \fill (3.80,1.45) circle (1.5pt) node[right=3pt] {$b$};
  \node[align=center,text width=1.55cm,fill=white,inner sep=1pt] at (3.45,0.78) {$u(a)=1$；\\$b\succ a$\\$\Rightarrow u(b)=2$};
  \node[align=center] at (3.45,0.25) {插入新元素};
  \draw[->,thick] (4.48,1.3) -- (4.92,1.3);
  \begin{scope}[xshift=8.0cm]
    \draw[thick,fill=gray!15] (0,0) rectangle (2.3,2.6);
    \node[font=\small\bfseries] at (1.15,2.28) {$n=3$};
    \fill (0.80,1.45) circle (1.5pt) node[left=3pt] {$a$};
    \fill (1.50,1.45) circle (1.5pt) node[below=3pt] {$c$};
    \fill (2.20,1.45) circle (1.5pt) node[right=3pt] {$b$};
    \node[align=center] at (1.15,0.75) {$c\succ a$，$c\sim b$\\$\Rightarrow u=(1,2,2)$};
    \node[align=center] at (1.15,0.25) {重新标号，保持保序};
  \end{scope}
\end{tikzpicture}
\caption{有限集效用表示定理：归纳构造的三个阶段}
\end{figure}
\figexplan{三个面板展示归纳构造的递推：$n=1$ 时取 $u(a)=1$；$n=2$ 时把新元素 $b$ 与已有元素比较，若 $b\succ a$ 则赋 $u(b)=2$；$n=3$ 时新元素 $c$ 满足 $c\succ a$ 且 $c\sim b$，于是得到 $u=(1,2,2)$。关键在于每一步只对“新元素与已编号元素”作比较，再重标号使保序成立——这正是有限集上表示定理的构造性证明。}

\subsection{存在性定理与两种构造证明}

\begin{theorem}[有限集表示定理]\label{thm:finite}
设 $X$ 为有限集，$\succeq$ 为 $X$ 上的弱序（完备 + 传递）。
则存在 $u:X\to\R$ 表示 $\succeq$。
\end{theorem}

给出两个证明：\textbf{归纳插入法}（揭示无差异类的层结构）与
\textbf{排名计数法}（给出显式公式）。两者都产出具整数取值的效用函数。

\begin{proof}[证明一：归纳插入法]
对 $|X|=n$ 作数学归纳。

\textbf{基例} $n=1$：$X=\{x_1\}$，取 $u(x_1)=0$ 即可。

\textbf{归纳步}：设定理对 $n-1$ 元集成立。任取 $X$ 中一个“最差”元：
有限弱序必有最小元（否则任取 $x^{(1)}$，若它非最小，存在 $x^{(2)}$ 使 $x^{(1)}\succ x^{(2)}$，
如此递降，有限性迫使过程终止），记为 $x_{\min}$。
令 $X'=X\setminus\{x_{\min}\}$，$\succeq$ 限制在 $X'$ 上仍是弱序。
由归纳假设，存在 $u':X'\to\R$ 表示 $\succeq|_{X'}$。延拓：
\[
u(x)=
\begin{cases}
u'(x), & x\in X',\\[2pt]
\min_{z\in X'} u'(z)-1, & x=x_{\min}.
\end{cases}
\]
验证：在 $X'$ 内部比较不变；对 $z\in X'$，$z\succeq x_{\min}$ 且
$u(z)=u'(z)\ge \min u' > u'(x_{\min})$ 方向的数值比较成立；
若 $z\sim x_{\min}$ 则把 $x_{\min}$ 改赋 $\min u'$ 即可（同值对应无差异）。
严格情形 $z\succ x_{\min}\iff u(z)>u(x_{\min})$ 由构造直接成立。故 $u$ 表示 $\succeq$。
\end{proof}

\begin{proof}[证明二：排名计数法]
定义\textbf{下截面计数函数}
\[
u(x)\ \coloneqq\ \#\{y\in X:\ x\succeq y\},
\]
即“被 $x$ 击败（含打平）的方案数”。验证它表示 $\succeq$。

\textbf{(i) $x\sim y\Rightarrow u(x)=u(y)$}：
$x\succeq z\iff y\succeq z$（$x\sim y$ 时 $x\succeq z\Rightarrow y\succeq x\succeq z\Rightarrow y\succeq z$，反向对称），
故两计数集相同，计数相等。

\textbf{(ii) $x\succ y\Rightarrow u(x)>u(y)$}：
记 $A\coloneqq\{z:y\succeq z\}$、$B\coloneqq\{z:x\succeq z\}$。
由传递性，$y\succeq z\Rightarrow x\succeq y\succeq z\Rightarrow x\succeq z$，故 $A\subseteq B$。
又 $x\in B$（自反性 $x\succeq x$）但 $x\notin A$（若 $y\succeq x$，与 $x\succ y$ 矛盾），
故 $A\subsetneq B$，两个有限集计数满足 $u(x)=|B|>|A|=u(y)$。

\textbf{(iii) $u(x)\ge u(y)\Rightarrow x\succeq y$}：
由完备性，若不然则 $y\succ x$，由 (ii) 得 $u(y)>u(x)$，矛盾。

三步合起来：$x\succeq y\iff u(x)\ge u(y)$。
\end{proof}

\begin{remark}
证明一给出的效用值是\textbf{逐层插入}的整数（同一无差异类同层同值）；
证明二的计数函数在同一无差异类上自动同值（步骤 (i)），
在严格偏好对之间拉开\textbf{整数间隔}（步骤 (ii) 的严格包含）。
两者都表明：有限集上不仅存在性成立，而且表示可以取为
$u:X\to\{0,1,\dots,n-1\}$ 的“名次函数”。
\end{remark}

\subsection{数值例：四种饮料的偏好}

\begin{example}\label{ex:finite}
设 $X=\{a,b,c,d\}$，偏好为 $a\succ b\sim c\succ d$（$a$ 严格最优，$b,c$ 无差异，$d$ 最差）。
\begin{itemize}
  \item \textbf{归纳构造}：最差元 $d$ 赋 $0$；$b\sim c$ 次之，同赋 $1$；$a$ 最优，赋 $2$：
        \[
        u(a)=2,\quad u(b)=u(c)=1,\quad u(d)=0.
        \]
  \item \textbf{排名计数}：$u(a)=\#\{a,b,c,d\}=4$，$u(b)=u(c)=\#\{b,c,d\}=3$，$u(d)=\#\{d\}=1$。
\end{itemize}
两者都满足表示定义；它们之间由严格递增映射 $f(0)=1,f(1)=3,f(2)=4$ 联系（命题~\ref{prop:ordinal}）。
\end{example}

\begin{center}
\begin{tikzpicture}[scale=1.16,transform shape]
  \draw[axisline] (0,0) -- (0,3.4) node[above, axislabel] {效用值};
  \foreach \y/\lab in {0.5/0, 1.5/1, 2.5/2}{
    \draw[AxisBlue] (0.08,\y) -- (-0.08,\y) node[left, axislabel] {\lab};
  }
  \foreach \x/\name/\val in {1.2/a/2.5, 2.6/b/1.5, 4.0/c/1.5, 5.4/d/0.5}{
    \node[nodeblue, label=above:{\name}] at (\x,\val) {};
  }
  \draw[dashed, MidBlue] (1.2,2.5) -- (0,2.5);
  \draw[dashed, MidBlue] (2.6,1.5) -- (0,1.5);
  \draw[dashed, MidBlue] (4.0,1.5) -- (0,1.5);
  \draw[dashed, MidBlue] (5.4,0.5) -- (0,0.5);
  \node at (3.2,-0.7) {$b$ 与 $c$ 同层（无差异类 $\Rightarrow$ 同效用）；层间间隔任意，只需保持顺序};
\end{tikzpicture}
\end{center}
\figexplan{纵轴为效用值。四个备选 $a,b,c,d$ 按偏好高低落在不同层上：$a$ 最高，$b$ 与 $c$ 同层（同一无差异类对应同一效用值），$d$ 最低。虚线把每个点投影到纵轴读数。层内点的具体数值与层间间隔均任意，只需保持偏好顺序——这正是序数性的含义。}

\subsection{有限集结论的边界}

\begin{remark}[三个容易误推的扩展]\label{rem:finite-boundary}
\textbf{(i) 无限集上归纳法失效}：证明一的“取最差元”在无限集上可能失败——
如 $\R_+$ 上的通常大小关系没有最小元；可数集 $(0,1]$ 上也没有最小元。
无限情形需要全新的构造（§3.3）。

\textbf{(ii) 不传递时结论崩溃}：设 $X=\{a,b,c\}$，$a\succ b\succ c\succ a$（循环）。
任何表示都要求 $u(a)>u(b)>u(c)>u(a)$，不可能。
这从构造层面复核了命题~\ref{prop:rep-rational}：传递性不可缺。

\textbf{(iii) 不完备时排名计数失真}：若 $a$ 与 $b$ 不可比，
计数法无法区分“无差异”与“不可比”——而效用函数强制把一切编码为实数比较，
不可比对 $(a,b)$ 必然被错误地归类为 $u(a)\ge u(b)$ 或反之。完备性同样不可缺。
\end{remark}

% ============================================================
\section{无限集：偏好仍可能被效用表示}
% ============================================================

有限集证明的两条腿——“取最差元”与“有限计数”——在无限集上都失效。
但\textbf{无限本身不构成障碍}：$\R_+$ 上的通常大小关系由 $u(x)=x$ 平凡表示。
真正的问题是：哪些无限结构允许数值化？本节按基数从小到大推进。

\subsection{可数无限集：权数构造}

\begin{theorem}[可数集表示定理]\label{thm:countable}
设 $X$ 为至多可数集，$\succeq$ 为 $X$ 上的弱序。则 $\succeq$ 可效用表示。
\end{theorem}

\begin{proof}
设 $X=\{x_1,x_2,x_3,\dots\}$ 为一个枚举（有限情形已由定理~\ref{thm:finite} 覆盖）。
给第 $n$ 个元素分配\textbf{几何权数} $w_n=2^{-n}$，定义
\[
u(x)\ \coloneqq\ \sum_{n:\,x\succ x_n} 2^{-n}\ -\ \sum_{n:\,x_n\succ x} 2^{-n},
\]
即“被 $x$ 击败者的权数和”减去“击败 $x$ 者的权数和”。两级数正项且总权 $\sum 2^{-n}\le 1$，
故 $u(x)$ 是良定义的实数。验证表示性。

\textbf{(i) $x\succ y\Rightarrow u(x)>u(y)$}：定义差集指标
\[
S(x,y)\ \coloneqq\ \{n:\ x\succeq x_n\ \text{且}\ y\not\succeq x_n\}\ \cup\ \{n:\ x_n\succeq y\ \text{且}\ x_n\not\succeq x\}.
\]
由传递性，$x\succ y$ 时“$x$ 的下截面”严格包含“$y$ 的下截面”、
“$y$ 的上截面”严格包含“$x$ 的上截面”，故 $u(x)-u(y)$ 是两个正项级数之差，
且每个指标 $n$ 对差值的贡献非负、至少一个指标贡献严格为正（$y$ 自身满足 $x\succ y$，
对应指标贡献 $2^{-n(y)}$ 于第一部分）。于是
\[
u(x)-u(y)\ \ge\ 2^{-n(y)}\ >\ 0.
\]

\textbf{(ii) $u(x)>u(y)\Rightarrow x\succ y$}：由完备性，若 $y\succeq x$，
则 $x\succ y$ 不成立，且分两种情形：$x\sim y$ 时两者的上、下截面完全相同
（$x\sim y$ 与传递性给出 $x\succeq z\iff y\succeq z$），故 $u(x)=u(y)$，矛盾；
$y\succ x$ 时由 (i) 得 $u(y)>u(x)$，矛盾。完备性排除了第三种可能，故 $x\succ y$。

合起来 $x\succeq y\iff u(x)\ge u(y)$。
\end{proof}

\begin{remark}[权数构造的思想]
给每个元素一个“可求和的权重”，把\textbf{集合包含关系}翻译成\textbf{数值大小关系}。
几何权数 $2^{-n}$ 保证任意子集的权和有限，差值的正负完全由包含方向决定。
同一个思想将在 §3.4 升级为：先把一个可数序稠密集“保序编号”进 $(0,1)$，再延拓到整个 $X$。
\end{remark}

\begin{example}[无限阶梯]
$X=\N$，偏好 $1\prec 2\prec 3\prec\cdots$（越大越好）。
权数构造给出 $u(k)=\sum_{n<k}2^{-n}=1-2^{-(k-1)}$，取值于 $[0,1)$；
当然 $u(k)=k$ 也是合法表示——再次印证表示不唯一（命题~\ref{prop:ordinal}），
严格递增变换 $f(t)=-\ln(1-t)$（$t<1$）把前者映为后者的平移。
\end{example}

\subsection{连续统：连续偏好与连续效用}

对 $X=\R^n_+$ 这类不可数消费集，权数构造无法直接套用（不可数正项和发散）。
经济学中的标准出路是\textbf{拓扑条件}：用连续性替代可数性。

\begin{definition}[连续偏好]\label{def:cont-pref}
设 $X\subseteq\R^n$。弱序 $\succeq$ \textbf{连续}，当且仅当对每个 $x\in X$，
上轮廓集 $\{y\in X:y\succeq x\}$ 与下轮廓集 $\{y\in X:x\succeq y\}$ 都是 $X$ 中的闭集。
等价地：若 $y^k\succeq x$（或 $x\succeq y^k$）且 $y^k\to y$，则 $y\succeq x$（或 $x\succeq y$）。
\end{definition}

\begin{remark}[连续性的序含义]
“极限运算不翻转偏好方向”：序列中每一项都至少和 $x$ 一样好，极限点不会突然变得比 $x$ 差。
词典序偏好正是破坏这一点的典型。取
\[
y^k=\Bigl(\frac{1}{k},\,1\Bigr),\qquad x=(0,\,2).
\]
对每个 $k$，$y^k$ 的第一坐标 $1/k>0$，故 $y^k\succ_L x$，即 $y^k$ 全落在 $x$ 的上轮廓集
$U(x)=\{z:z\succeq_L x\}$ 中；但 $y^k\to(0,1)$，而 $(0,1)\prec_L(0,2)=x$，极限点\textbf{掉出}了 $U(x)$。
故 $U(x)$ 不是闭集，$\succeq_L$ 不连续。
这从拓扑角度重新解释了定理~\ref{thm:lex}：不连续的偏好不可能有连续效用表示；
而 §\ref{subsec:lex} 的基数论证更强——它连不连续的效用表示都没有。
\end{remark}

\begin{theorem}[连续偏好的效用表示定理]\label{thm:continuous-rep}
设 $X\subseteq\R^n$ 为非空闭凸集（或更一般地，连通可分拓扑空间），
$\succeq$ 为 $X$ 上连续的弱序。则存在\textbf{连续}效用函数 $u:X\to\R$ 表示 $\succeq$。
\end{theorem}

证明分三步，核心步骤（连续编号）在 §3.4 完整展开：

\begin{enumerate}
  \item \textbf{可数序稠密子集的存在性}：$X$ 可分（含可数稠密集 $D$），
        连续性把“$D$ 稠密”升级为“$D$ 序稠密”：
        若 $x\succ y$，连续性给出开球 $B(x,\varepsilon)$ 内所有点仍 $\succ y$，取 $d\in D\cap B(x,\varepsilon)$，
        则 $x\succeq d\succeq y$（严格说需多一步插值，见 §3.4 引理~\ref{lem:interp}）。
  \item \textbf{连续编号}：把 $D$ 保序地映成 $(0,1)$ 的稠密子集（§3.4 定义~\ref{def:cont-num}）。
  \item \textbf{延拓与连续性}：用 $u(x)=\sup\{g(d):d\in D,\ x\succeq d\}$ 延拓到 $X$；
        偏好连续 $\Rightarrow$ $u$ 的上下轮廓集闭 $\Rightarrow$ $u$ 连续（§3.4 定理~\ref{thm:u-cont}）。
\end{enumerate}

\begin{remark}[与必要条件定理的关系]
定理~\ref{thm:dense-necessary} 说可数序稠密子集\textbf{必要}；
本定理的证明路线说明在拓扑配合下它\textbf{也充分}。
事实上去掉全部拓扑假设后，“弱序 + 可数序稠密子集”本身就是充要条件（§3.4 完成其充分性证明）——
连续性是“让 $X$ 自带一个现成的可数序稠密集”的经济学方便假设，
并且额外买回了效用函数的连续性（这对最优化理论至关重要：
连续效用 + 紧预算集 $\Rightarrow$ 最优解存在）。
\end{remark}

\subsection{单调性给出的显式表示 \texorpdfstring{$\star$}{*}}

对单调偏好，表示可以更具体。

\begin{proposition}[单调连续偏好的对角线构造]\label{prop:diagonal}
设 $\succeq$ 为 $\R^n_+$ 上连续、强单调的弱序。
则每个无差异类与对角线 $\{t\mathbf{1}:t\ge0\}$（$\mathbf{1}=(1,\dots,1)$）恰交于一点，
从而
\[
u(x)\ \coloneqq\ t\ \text{使得}\ x\sim t\mathbf{1}
\]
是表示 $\succeq$ 的连续效用函数。
\end{proposition}

\begin{proof}[证明要点]
\textbf{存在性}：设 $\bar t=\max_i x_i$。由强单调性 $\bar t\mathbf{1}\succeq x\succeq 0\cdot\mathbf{1}$。
考察 $g(t)=\mathbf{1}_{t\mathbf{1}\succeq x}$：连续性使 $\{t:t\mathbf{1}\succeq x\}$ 为闭区间
（含 $\bar t$，下确界记为 $t^*$），$\{t:x\succeq t\mathbf{1}\}$ 亦为闭区间（含 $0$，上确界记为 $t^{*}$）。
完备性迫使两闭区间覆盖 $[0,\bar t]$，交点处 $t\mathbf{1}\sim x$。
\textbf{唯一性}：$t_1<t_2\Rightarrow t_2\mathbf{1}\gg t_1\mathbf{1}\Rightarrow t_2\mathbf{1}\succ t_1\mathbf{1}$（强单调），
故无差异类至多交对角线一次。
\textbf{表示性}：$x\succeq y\iff u(x)\mathbf{1}\sim x\succeq y\sim u(y)\mathbf{1}\iff u(x)\mathbf{1}\succeq u(y)\mathbf{1}\iff u(x)\ge u(y)$，
末步用强单调性。
\textbf{连续性}：$u$ 的上、下轮廓集即 $\succeq$ 的上、下轮廓集，由连续性为闭集。
\end{proof}

\begin{center}
\begin{tikzpicture}[scale=1.15]
  \draw[axisline] (0,0) -- (5.2,0) node[below left, axislabel] {$x_1$};
  \draw[axisline] (0,0) -- (0,5.2) node[below left, axislabel] {$x_2$};
  % diagonal
  \draw[thick, dashed, AxisBlue] (0,0) -- (4.9,4.9) node[right, align=left, axislabel] {对角线\\ $t\mathbf{1}$};
  % indifference curves
  \draw[thick, MidBlue, domain=0.55:4.6, smooth, variable=\x] plot ({\x},{1.0/\x*2.4});
  \draw[thick, MidBlue, domain=0.85:4.6, smooth, variable=\x] plot ({\x},{1.0/\x*4.2});
  \draw[thick, MidBlue, domain=1.3:4.6, smooth, variable=\x] plot ({\x},{1.0/\x*6.4});
  % intersection points
  \node[nodeblue, label=below right:{$u(x')\mathbf{1}$}] at (1.549,1.549) {};
  \node[nodeblue, label=below right:{$u(x)\mathbf{1}$}] at (2.049,2.049) {};
  \node[nodeblue] at (2.53,2.53) {};
  % a bundle x
  \node[nodeblue, label=right:{$x$}] at (3.6,1.1667) {};
  \node[align=left] at (4.3,4.2) {无差异曲线与对角线\\ 恰交于一点（强单调）};
\end{tikzpicture}
\end{center}
\figexplan{对角线 $t\mathbf{1}$ 是构造效用值的标尺：强单调性保证每条无差异曲线与对角线恰交于一点，该交点 $u(x)\mathbf{1}$ 的数值即定义为 $x$ 的效用。图中三条双曲型无差异曲线由低到高排列，越靠右上效用越高；$x$ 所在无差异曲线与对角线的交点给出 $u(x)$。}

命题~\ref{prop:diagonal} 的构造把“效用值”定义为\textbf{等值的对角线消费束的坐标}，
这是福利分析中“等值变化”的雏形。同一思想在 §3.7 拟线性情形将精确化为
$u(t,y)=t+v(y)$，其中 $v(y)$ 就是“用计价物度量的 $y$ 的价值”。

% ============================================================
\section{连续编号：从可数稠密集到效用函数}\label{sec:numbering}
% ============================================================

本节完成定理~\ref{thm:dense-necessary} 的逆命题（充分性）与
定理~\ref{thm:continuous-rep} 的核心步骤。构造分两段：
先把可数序稠密集 $D$ \textbf{保序地}嵌入 $(0,1)$（连续编号），
再沿偏好序把编号\textbf{延拓}到全空间 $X$。

\subsection{连续编号的构造}

\begin{lemma}[插值引理]\label{lem:interp}
设 $\succeq$ 为弱序，$D$ 为 $X$ 的可数序稠密子集。
若 $x\succ y$ 且 $x,y$ 之间\textbf{不是}“相邻跳跃对”（即不存在 $x\succ y$ 使
$\{z:x\succ z\succ y\}=\varnothing$），则存在 $d\in D$ 使 $x\succ d\succ y$。
\end{lemma}

\begin{proof}
$x\succ y$ 且非相邻跳跃，故存在 $w$ 使 $x\succ w\succ y$。
由序稠密定义，存在 $d_1\in D$ 使 $x\succeq d_1\succeq w$；再对 $w\succ y$ 用序稠密，
存在 $d_2\in D$ 使 $w\succeq d_2\succeq y$。若 $d_1\succ w$ 或 $w\succ d_2$ 之一成立，
则 $d_1$ 或 $d_2$ 已满足严格夹逼；若 $d_1\sim w\sim d_2$，则 $d_1\sim d_2\sim w$，
$x\succ d_1\succ y$ 直接成立（$d_1\sim w$ 且 $x\succ w\succ y$ 加传递性）。
\end{proof}

\begin{definition}[连续编号]\label{def:cont-num}
设 $D=\{d_1,d_2,\dots\}$ 为 $X$ 的可数序稠密子集的枚举。
\textbf{连续编号}是如下归纳构造的映射 $g:D\to(0,1)$：

\textbf{第 1 步}：令 $g(d_1)=1/2$。

\textbf{第 $n$ 步}：设 $g$ 已在 $\{d_1,\dots,d_{n-1}\}$ 上定义且保序。
$d_n$ 与已编号元素的关系必属三种情形之一：
\begin{enumerate}
  \item \textbf{新最大}：$d_n\succ d_i$ 对所有 $i<n$。令
        $g(d_n)=\dfrac{1+\max_{i<n}g(d_i)}{2}$（取当前最大值与 $1$ 的中点）；
  \item \textbf{新最小}：$d_i\succ d_n$ 对所有 $i<n$。令
        $g(d_n)=\dfrac{0+\min_{i<n}g(d_i)}{2}$（取 $0$ 与当前最小值的中点）；
  \item \textbf{中间}：存在已编号元夹逼 $d_n$。取序相邻的一对
        $d_L\prec d_n\prec d_R$（$d_L,d_R\in\{d_1,\dots,d_{n-1}\}$ 且二者之间无其他已编号元），令
        \[
        g(d_n)=\frac{g(d_L)+g(d_R)}{2};
        \]
        若 $d_n$ 与某已编号元无差异（$d_n\sim d_i$），则令 $g(d_n)=g(d_i)$。
\end{enumerate}
\end{definition}

\begin{proposition}[连续编号保序]\label{prop:num-order}
定义~\ref{def:cont-num} 的 $g$ 满足：对任意 $d,d'\in D$，
\[
d\succeq d'\ \Longleftrightarrow\ g(d)\ge g(d'),\qquad
d\succ d'\ \Longleftrightarrow\ g(d)>g(d').
\]
且 $g(D)$ 在 $(0,1)$ 中稠密：$g(D)$ 与 $D$ 同序型地铺满——
任意 $g(d)<g(d')$ 之间存在其他编号值（不断二分产生）。
\end{proposition}

\begin{proof}
归纳：第 $n$ 步放入 $g(d_n)$ 时，三种情形的取值都严格落在由偏好序决定的
目标开区间内（新最大取 $(\max g,1)$ 中点、新最小取 $(0,\min g)$ 中点、
中间取 $(g(d_L),g(d_R))$ 中点），故插入后序关系保持。
无差异情形赋同值，故 $d\sim d'\iff g(d)=g(d')$。
稠密性：若 $d\succ d'$ 非相邻跳跃，由引理~\ref{lem:interp} 存在 $d''\in D$ 使
$d\succ d''\succ d'$，保序性给出 $g(d)>g(d'')>g(d')$。
\end{proof}

\begin{center}
\begin{tikzpicture}[scale=1.0]
  \draw[axisline] (0,0) -- (10.6,0) node[below left, axislabel] {$(0,1)$};
  \draw[AxisBlue] (0.3,0.08) -- (0.3,-0.08) node[below, axislabel] {$0$};
  \draw[AxisBlue] (10.3,0.08) -- (10.3,-0.08) node[below, axislabel] {$1$};
  % 下层标签
  \foreach \x/\lab in {5.3/{$g(d_1)=\frac12$}, 2.8/{$g(d_3)=\frac14$}}{
    \node[nodeblue, label=above:{\lab}] at (\x,0) {};
  }
  % 错层避免标签碰撞
  \node[nodeblue, label=above=14pt:{$g(d_4)=\frac58$}] at (6.55,0) {};
  % 上层标签（与下层错开）
  \foreach \x/\lab in {7.8/{$g(d_2)=\frac34$}, 4.05/{$g(d_5)=\frac38$}}{
    \node[nodeblue] at (\x,0) {};
    \node[above=15pt] at (\x,0) {\lab};
    \draw[thin, MidBlue] (\x,0.12) -- (\x,0.62);
  }
  \node[align=center] at (5.3,-1.1) {编号顺序：$d_1\to d_2\to d_3\to d_4\to d_5$；每一步把新元素放入\\ 偏好序指定的区间\textbf{中点}，编号值自动保序且越来越密};
\end{tikzpicture}
\end{center}
\figexplan{横轴为 $(0,1)$ 区间。可数稠密集的元素按偏好指定的顺序被逐个插入相应区间的中点，得到编号值 $g(d_i)$（如 $g(d_1)=\frac12$、$g(d_2)=\frac34$）。所有点都画在同一条轴上，标签上下错层仅为避免文字碰撞，并不表示不同含义；随着插入进行，编号自动保序且越来越密。}

\subsection{延拓：从 $D$ 到 $X$}

\begin{theorem}[延拓表示定理]\label{thm:extension}
设 $\succeq$ 为弱序且 $X$ 含可数序稠密子集。则 $\succeq$ 可效用表示。
\end{theorem}

\begin{proof}
分三步：补足稠密集、连续编号、上确界延拓。

\textbf{第 1 步（把缺口端点并入稠密集）。}
先证：对任意可数序稠密集 $D$，$\succeq$ 的相邻跳跃对（定义~\ref{def:gap}）至多可数。
事实上，每个相邻跳跃对 $x\succ y$ 由序稠密性给出 $d\in D$ 使 $x\succeq d\succeq y$；
相邻性迫使 $d\sim x$ 或 $d\sim y$，即 $d$ 落在跳跃的某一端。
而同一 $d$ 至多充当两个跳跃的端点（上方一个、下方一个），
不同跳跃对必须占用不同的 $d$，故跳跃对数 $\le 2|D|$，至多可数。
从每个跳跃对的\textbf{两端}无差异类各取一个代表元，并入 $D$，得到新的可数集 $D'$。
$D'$ 仍序稠密，且\textbf{每个跳跃对的两端都有 $D'$ 点}。

\textbf{第 2 步（连续编号）。}
对 $D'$ 作定义~\ref{def:cont-num} 的连续编号 $g:D'\to(0,1)$，
由命题~\ref{prop:num-order}，$g$ 严格保序：$d\succ d'\iff g(d)>g(d')$。

\textbf{第 3 步（上确界延拓）。}
定义
\[
u(x)\ \coloneqq\ \sup\,\{g(d):\ d\in D',\ x\succeq d\},
\qquad \sup\varnothing\coloneqq 0.
\]
验证表示性：

\textbf{(i) $x\succeq y\Rightarrow u(x)\ge u(y)$}：
传递性给出 $\{d:y\succeq d\}\subseteq\{d:x\succeq d\}$，上确界对包含单调。

\textbf{(ii) $x\succ y\Rightarrow u(x)>u(y)$}：分两种情形。

\textit{情形 1：存在 $z$ 使 $x\succ z\succ y$（非相邻）。}
由序稠密性取 $d_1\in D'$ 使 $x\succeq d_1\succeq z$，则 $d_1\succeq z\succ y$。
再对 $d_1\succ y$ 用一次：若 $d_1\succ y$ 之间有 $z'$，取 $d_2\in D'$ 使
$d_1\succeq z'\succeq d_2\succ y$ 的夹逼形式 $d_1\succ d_2\succeq y$ 或 $d_1\succeq d_2\succ y$；
无论哪种，都有 $D'$ 中的两点 $d_1\succ d_2$ 且 $x\succeq d_1$、$d_2\succ y$。
若 $d_1\succ y$ 是相邻跳跃，第 1 步已保证两端都有 $D'$ 点，
取 $d_2\in D'$、$d_2\sim y$，同样有 $d_1\succ d_2$。
统一地：存在 $d_1,d_2\in D'$，$x\succeq d_1\succ d_2\succeq y$。于是
\[
u(x)\ \ge\ g(d_1)\ >\ g(d_2)\ \ge\ u(y),
\]
末步因任何 $y\succeq d$ 都有 $d_2\succeq y\succeq d$ 即 $g(d_2)\ge g(d)$，
故 $g(d_2)$ 是 $\{g(d):y\succeq d\}$ 的上界。

\textit{情形 2：$x\succ y$ 为相邻跳跃对。}
第 1 步保证存在 $d_x,d_y\in D'$ 使 $d_x\sim x$、$d_y\sim y$。则
\[
u(x)\ \ge\ g(d_x)\ >\ g(d_y)\ =\ u(y),
\]
其中等号 $u(y)=g(d_y)$：任何 $y\succeq d$ 都有 $d\preceq y$
（否则 $d\succ y$ 且相邻性迫使 $d\succeq x$，即 $d\succeq x\succ y$，与 $y\succeq d$ 矛盾），
故 $g(d)\le g(d_y)$，上确界为 $g(d_y)$。

\textbf{(iii) $u(x)\ge u(y)\Rightarrow x\succeq y$}：
若不然，完备性给出 $y\succ x$，由 (ii) 得 $u(y)>u(x)$，矛盾。
\end{proof}

\subsection{缺口分析：效用函数的跳跃从何而来}

连续编号产出的 $u$ 一般是\textbf{有跳跃的}阶梯型函数。跳跃不是构造的缺陷，
而是偏好序本身的“缺口”的忠实影像。

\begin{definition}[序缺口]\label{def:gap}
弱序 $\succeq$ 的一个\textbf{缺口}是一对无差异类 $([y],[x])$ 满足
$x\succ y$ 且 $\{z:x\succ z\succ y\}=\varnothing$（相邻跳跃对），
或 $(-\infty,[x])$、$([y],+\infty)$ 型的首、尾端点。
\end{definition}

\begin{proposition}[缺口至多可数]\label{prop:gap-countable}
若 $\succeq$ 可表示，则其缺口族至多可数。
\end{proposition}

\begin{proof}
每个缺口 $([y],[x])$ 对应实数轴上的非空开区间 $\bigl(u(y),u(x)\bigr)$
（$u$ 为任一表示），不同缺口对应的开区间两两不交
（缺口之间由序的相邻性决定端点，端点不重叠）。
实数轴上两两不交的非空开区间族至多可数（每区间取一有理数即得单射入 $\Q$）。
\end{proof}

\begin{center}
\begin{tikzpicture}[scale=1.0]
  \draw[axisline] (0,0) -- (0,3.6) node[above, axislabel] {$u(x)$};
  \draw[axisline] (0,0) -- (8.8,0) node[below left, axislabel] {$X$（按 $\succeq$ 排序）};
  % steps with jumps
  \draw[thick, AxisBlue] (0.6,0.5) -- (2.2,0.5);
  \draw[thick, dashed, AxisBlue] (2.2,0.5) -- (2.2,1.4);
  \draw[thick, AxisBlue] (2.2,1.4) -- (4.0,1.4);
  \draw[thick, dashed, AxisBlue] (4.0,1.4) -- (4.0,2.6);
  \draw[thick, AxisBlue] (4.0,2.6) -- (6.2,2.6);
  \draw[thick, dashed, AxisBlue] (6.2,2.6) -- (6.2,3.1);
  \draw[thick, AxisBlue] (6.2,3.1) -- (8.2,3.1);
  \node[nodeblue] at (2.2,0.5) {}; \draw[DeepBlue, fill=white, thick] (2.2,1.4) circle (1.8pt);
  \node[nodeblue] at (4.0,1.4) {}; \draw[DeepBlue, fill=white, thick] (4.0,2.6) circle (1.8pt);
  \node[nodeblue] at (6.2,2.6) {}; \draw[DeepBlue, fill=white, thick] (6.2,3.1) circle (1.8pt);
  % jump intervals on y-axis
  \draw[decorate, decoration={brace, amplitude=4pt, mirror}] (-0.15,0.5) -- (-0.15,1.4) node[midway, left=6pt, align=center] {缺口 1\\ $(0.5,1.4)$};
  \draw[decorate, decoration={brace, amplitude=4pt, mirror}] (-0.15,1.4) -- (-0.15,2.6) node[midway, left=6pt, align=center] {缺口 2};
  \node[align=center] at (4.4,-0.9) {每个跳跃对应值域轴上一个非空开区间；区间两两不交 $\Rightarrow$ 至多可数个跳跃};
\end{tikzpicture}
\end{center}
\figexplan{阶梯折线是效用函数沿偏好序的图象：水平平台段表示同一无差异类（同一效用值），竖直虚线段表示跳跃——相邻层级之间没有中间元素，对应值域轴上一个非空开区间（左侧花括号标出的“缺口”）。实心点取到该平台值、空心点表示跳跃另一端未被取到；两两不交的开区间至多可数，故跳跃至多可数个。}

\begin{theorem}[连续性恢复]\label{thm:u-cont}
设 $X\subseteq\R^n$ 连通，$\succeq$ 为连续弱序。则 $\succeq$ 没有缺口，
且上确界延拓给出的表示是\textbf{连续}函数。
\end{theorem}

\begin{proof}[证明要点]
\textbf{无缺口}：设 $x\succ y$。连续性使 $V=\{z:x\succ z\succ y\}$ 为
两个开集之交，从而是开集；它非空——连通集上的连续弱序不存在相邻跳跃对
（否则 $X=\{z:z\succeq x\}\cup\{z:y\succeq z\}$ 把 $X$ 分成两个非空闭集且
$\{z:z\succeq x\}\cap\{z:y\succeq z\}=\{z:z\sim x\ \text{或}\ z\sim y\}$ 的边界行为与连通性矛盾；
标准论证：开集 $\{z:x\succ z\}$ 与 $\{z:z\succ y\}$ 的真包含链给出中间点）。
故任意 $x\succ y$ 之间都有 $z$，无缺口。

\textbf{延拓连续}：无缺口时，对任意实数 $r\in(0,1)$，
$\{x:u(x)\ge r\}=\{x:x\succeq z\ \text{对所有}\ u(z)<r\ \text{的}\ z\}$
是闭集（连续偏好的上轮廓集之交）；
$\{x:u(x)\le r\}$ 对偶地为闭集。$u$ 的上、下水平集皆闭，故 $u$ 连续。
\end{proof}

\begin{remark}[连续编号法的全貌]
\[
\underbrace{\text{弱序}+\text{可数序稠密集}}_{\text{纯序条件}}
\xrightarrow[\text{中点插入}]{\text{连续编号 }g}
\underbrace{g(D)\subset(0,1)\ \text{稠密}}_{\text{保序}}
\xrightarrow[\sup\text{ 延拓}]{}
\underbrace{u:X\to[0,1]\ \text{表示}}_{\text{存在性}}
\xrightarrow[\text{无缺口}]{\text{偏好连续}}
\underbrace{u\ \text{连续}}_{\text{可分析性}}.
\]
整条链条只用到序结构与“不交开区间至多可数”这一个实数性质，
是效用理论中\textbf{假设最省}的构造路线。
\end{remark}

% ============================================================
\section{偏好的特性如何翻译为效用函数的性质}\label{sec:properties}
% ============================================================

设 $X=\R^n_+$（或 $\R^n_{++}$），$\succeq$ 为连续弱序，$u$ 为其连续表示。
本节逐条建立“偏好语言 $\leftrightarrow$ 效用语言”的翻译词典。
每一条都回答两个问题：\textbf{充要条件是什么？}、\textbf{该性质是序数的还是基数的？}

\subsection{单调性 $\leftrightarrow$ 递增性}

\begin{definition}[偏好的单调性]\label{def:monotone}
$\succeq$ \textbf{单调}：$x\gg y\Rightarrow x\succ y$，且 $x\ge y\Rightarrow x\succeq y$；
\textbf{强单调}：$x\ge y$ 且 $x\ne y\Rightarrow x\succ y$；
\textbf{局部非饱和}：$\forall x,\forall\varepsilon>0,\exists y\in B(x,\varepsilon):y\succ x$
（任何点的任意小邻域内都有更好点——排除“幸福点”与粗无差异带）。
三者关系：强单调 $\Rightarrow$ 单调 $\Rightarrow$ 局部非饱和。
\end{definition}

\begin{proposition}\label{prop:mono-u}
设 $u$ 连续表示 $\succeq$。则
\[
\succeq\ \text{单调}\ \Longleftrightarrow\ u\ \text{递增}:\ x\gg y\Rightarrow u(x)>u(y),\ x\ge y\Rightarrow u(x)\ge u(y);
\]
\[
\succeq\ \text{强单调}\ \Longleftrightarrow\ u\ \text{严格递增}:\ x\ge y,\ x\ne y\Rightarrow u(x)>u(y).
\]
\end{proposition}

\begin{proof}
由表示定义，$x\gg y$ 时 $x\succ y\iff u(x)>u(y)$；其余情形逐条同义替换即得。
\end{proof}

\begin{example}
$u(x_1,x_2)=x_1^{1/2}x_2^{1/2}$ 在 $\R^2_{++}$ 上严格递增
（$\partial u/\partial x_i>0$），对应强单调偏好；
$u(x_1,x_2)=\min\{x_1,x_2\}$ 递增但不严格递增（$x_1>x_2$ 时增大 $x_1$ 不改变效用），
对应单调但非强单调的偏好——$(2,1)$ 与 $(3,1)$ 无差异。
\end{example}

\begin{center}
\begin{tikzpicture}[scale=1.05]
  \draw[axisline] (0,0) -- (5.0,0) node[below left, axislabel] {$x_1$};
  \draw[axisline] (0,0) -- (0,4.6) node[below left, axislabel] {$x_2$};
  \node[nodeblue, label=below left:{$x$}] at (2.2,1.8) {};
  % better set (NE quadrant)
  \draw[fillsoft] (2.2,1.8) rectangle (4.8,4.4);
  \draw[dashed, MidBlue] (2.2,1.8) -- (2.2,4.4); \draw[dashed, MidBlue] (2.2,1.8) -- (4.8,1.8);
  \node[align=center] at (3.5,3.3) {$y\gg x$\\ $\Rightarrow u(y)>u(x)$};
  \draw[arrblue] (2.2,1.8) -- (3.1,2.7);
  \node[align=center] at (1.05,3.6) {单调性：上轮廓集\\ 包含东北象限};
\end{tikzpicture}
\end{center}
\figexplan{以 $x$ 为参照，浅灰矩形是东北方向的严格更优区域：任何 $y\gg x$ 都落在其中，单调性要求 $u(y)>u(x)$。两条虚线是过 $x$ 的横纵参考线，界定该区域边界；箭头指示由 $x$ 指向更优束的改进方向。注意 $\min$ 型效用递增但不严格递增，对应单调而非强单调。}

\subsection{凸性 $\leftrightarrow$ 拟凹性}

\begin{definition}[凸偏好]\label{def:convex-pref}
$\succeq$ \textbf{凸}：对任意 $x$ 与 $\lambda\in[0,1]$，
$y\succeq x,\ z\succeq x\ \Rightarrow\ \lambda y+(1-\lambda)z\succeq x$；
\textbf{严格凸}：$y\succeq x,\ z\succeq x,\ y\ne z,\ \lambda\in(0,1)\ \Rightarrow\ \lambda y+(1-\lambda)z\succ x$。
\end{definition}

\begin{definition}[拟凹]\label{def:quasiconcave}
$f:\R^n_+\to\R$ \textbf{拟凹}：$f(\lambda y+(1-\lambda)z)\ge\min\{f(y),f(z)\}$；
\textbf{严格拟凹}：$y\ne z,\lambda\in(0,1)$ 时不等式严格（当 $f(y)\ne f(z)$ 时按 $\min$ 端定义严格）。
\end{definition}

\begin{proposition}\label{prop:convex-qc}
$u$ 表示 $\succeq$。则
\[
\succeq\ \text{凸}\ \Longleftrightarrow\ u\ \text{拟凹};\qquad
\succeq\ \text{严格凸}\ \Longleftrightarrow\ u\ \text{严格拟凹}.
\]
\end{proposition}

\begin{proof}
\textbf{($\Rightarrow$)} 设 $y,z$ 任意，记 $m=\min\{f(y),f(z)\}$，不妨 $f(z)=m$，即 $u(y)\ge u(z)$ 也就是
$y\succeq z$。凸性给出 $\lambda y+(1-\lambda)z\succeq z$，
即 $u(\lambda y+(1-\lambda)z)\ge u(z)=m$。
\textbf{($\Leftarrow$)} 设 $y,z\succeq x$，则 $u(y),u(z)\ge u(x)$，
拟凹给出 $u(\lambda y+(1-\lambda)z)\ge\min\{u(y),u(z)\}\ge u(x)$，即凸偏好。
严格情形把 $\ge$ 全部换为 $>$（$\lambda\in(0,1)$、$y\ne z$）。
\end{proof}

\begin{center}
\begin{tikzpicture}[scale=1.05]
  \draw[axisline] (0,0) -- (5.4,0) node[below left, axislabel] {$x_1$};
  \draw[axisline] (0,0) -- (0,4.8) node[below left, axislabel] {$x_2$};
  % convex IC
  \draw[thick, MidBlue, domain=0.7:5.0, smooth, variable=\x] plot ({\x},{2.9/\x});
  \node at (4.55,1.15) {$u(x)=c$};
  \node[nodeblue, label=above left:{$y$}] at (1.45,2.0) {};
  \node[nodeblue, label=below right:{$z$}] at (4.0,0.725) {};
  \draw[dashed, MidBlue] (1.45,2.0) -- (4.0,0.725);
  \node[nodeblue, label=above right:{$\frac{y+z}{2}$}] at (2.725,1.3625) {};
  \node[align=left] at (3.9,3.6) {凸偏好：均值点落在\\ 无差异曲线\textbf{上方}\\ $\Rightarrow u\Bigl(\frac{y+z}{2}\Bigr)\ge\min\{u(y),u(z)\}$};
\end{tikzpicture}
\end{center}
\figexplan{双曲线为无差异曲线 $u(x)=c$，$y$ 与 $z$ 是其上两点，虚线是连接二者的混合线段。凸偏好要求线段中点 $\frac{y+z}{2}$ 落在无差异曲线上方，即 $u\bigl(\frac{y+z}{2}\bigr)\ge\min\{u(y),u(z)\}$——这正是效用函数拟凹性的几何表达。}

\subsection{位似性 $\leftrightarrow$ 齐次函数的单调变换}

\begin{definition}[位似偏好]\label{def:homothetic}
$\succeq$ \textbf{位似}（homothetic）：$x\sim y\ \Rightarrow\ \lambda x\sim\lambda y$ 对一切 $\lambda>0$。
等价地（结合单调性）：$x\succeq y\iff\lambda x\succeq\lambda y$。
\end{definition}

\begin{proposition}[位似偏好的表示特征]\label{prop:homothetic-u}
设 $\succeq$ 为 $\R^n_+$ 上连续、单调的弱序。则
\[
\succeq\ \text{位似}\ \Longleftrightarrow\ \succeq\ \text{可由一个一次齐次函数表示}
\ \Longleftrightarrow\ u=f\circ h,\ h\ \text{一次齐次},\ f\ \text{严格递增}.
\]
\end{proposition}

\begin{proof}
\textbf{($\Leftarrow$)} 若 $h$ 一次齐次表示 $\succeq$：
$x\sim y\Rightarrow h(x)=h(y)\Rightarrow h(\lambda x)=\lambda h(x)=\lambda h(y)=h(\lambda y)\Rightarrow\lambda x\sim\lambda y$。
单调变换 $f$ 不改变无差异关系（命题~\ref{prop:ordinal}）。

\textbf{($\Rightarrow$)} 用对角线构造（命题~\ref{prop:diagonal}）：
$u(x)$ 定义为使 $x\sim u(x)\mathbf{1}$ 的数。验证一次齐次性：
\[
\lambda x\sim\lambda\bigl(u(x)\mathbf{1}\bigr)=\bigl(\lambda u(x)\bigr)\mathbf{1}
\quad(\text{位似性})
\ \Longrightarrow\ u(\lambda x)=\lambda u(x),
\]
末步因对角线交点唯一。故对角线效用本身就是一次齐次表示。
\end{proof}

\begin{remark}[位似性的经济含义：射线上的边际替代率不变]
位似偏好的无差异曲线互为\textbf{径向放大}：沿从原点出发的任何射线，
无差异曲线的斜率（边际替代率）恒定。这等价于：收入扩张路径是过原点的直线，
需求收入弹性处处为 $1$。检验：$u(x_1,x_2)=x_1^{1/2}x_2^{1/2}$ 一次齐次（位似）；
$u(x_1,x_2)=x_1+\sqrt{x_2}$ 不是（§3.7 拟线性，收入扩张路径弯曲）。
\end{remark}

\begin{center}
\begin{tikzpicture}[scale=1.0]
  \draw[axisline] (0,0) -- (5.4,0) node[below left, axislabel] {$x_1$};
  \draw[axisline] (0,0) -- (0,4.8) node[below left, axislabel] {$x_2$};
  \draw[dashed, MidBlue] (0,0) -- (4.6,3.3) node[right, axislabel] {射线};
  \draw[thick, MidBlue, domain=0.8:5.0, smooth, variable=\x] plot ({\x},{2.4/\x});
  \draw[thick, MidBlue, domain=1.6:5.0, smooth, variable=\x] plot ({\x},{4.8/\x});
  % tangency slope marks along the ray
  \node[nodeblue] at (1.697,1.2129) {};
  \node[nodeblue] at (3.394,2.4258) {};
  \draw[AxisBlue] (1.197,1.7) -- (2.197,0.72);
  \draw[AxisBlue] (2.894,2.9) -- (3.894,1.94);
  \node[align=left] at (4.3,4.1) {同一射线上两处\\ 切线斜率相同\\（MRS 沿射线不变）};
\end{tikzpicture}
\end{center}
\figexplan{自原点出发的虚线射线依次穿过两条无差异曲线；两个交点处的切线（短实线）斜率相同。这表明沿任一射线边际替代率 MRS 保持不变，无差异曲线族互为沿射线的等比放大，即偏好位似；等价地，效用函数可写成齐次函数的单调变换。}

\subsection{可加可分性 $\leftrightarrow$ 加和形式（详证见 \S 3.6）}

\begin{definition}[可加可分表示]\label{def:additive-sep}
$\succeq$ \textbf{可加可分}：存在 $u_i:\R_+\to\R$（$i=1,\dots,n$）使
$u(x)=\sum_{i=1}^n u_i(x_i)$ 表示 $\succeq$。
\end{definition}

偏好侧的特征是\textbf{坐标独立性}：任意坐标子集上的偏好比较不受其余坐标取值影响。
$u(x_1,x_2)=\sqrt{x_1}+\sqrt{x_2}$ 可分；$u(x_1,x_2)=x_1x_2$ 不可分
（$x_1$ 的边际效用 $x_2$ 依赖于 $x_2$ 的水平——交叉偏导非零是破坏可分的可微判据）。
完整陈述与证明见 §3.6。

\subsection{拟线性 $\leftrightarrow$ 线性计价物结构（详证见 \S 3.7）}

\begin{definition}[拟线性表示]\label{def:ql-rep}
$\succeq$ \textbf{拟线性}：存在 $v:\R_+^{n-1}\to\R$ 使
$u(x_1,x_{-1})=x_1+v(x_{-1})$ 表示 $\succeq$，其中 $x_1$ 为计价物。
\end{definition}

偏好侧的特征是\textbf{无财富效应}：沿计价物方向平移不改变比较结果。
$u(x_1,x_2)=x_1+\sqrt{x_2}$ 拟线性；其无差异曲线互为\textbf{纵向平移}。
完整陈述与证明见 §3.7。

\subsection{连续性 $\leftrightarrow$ 连续表示（回顾）}

定义~\ref{def:cont-pref} 的连续偏好 $\iff$ 存在连续表示
（定理~\ref{thm:continuous-rep} 与定理~\ref{thm:u-cont}）。
可微性是进一步的\textbf{基数}便利：取到可微表示后，
边际替代率有公式
\[
\mathrm{MRS}_{ij}(x)=\frac{\partial u/\partial x_i}{\partial u/\partial x_j},
\]
该比值在严格递增变换下不变（链式法则中 $f'$ 约去），故 MRS 本身是\textbf{序数量}。

\subsection{序数性质与基数性质的严格区分 \texorpdfstring{$\star$}{*}}\label{subsec:ordinal-cardinal}

\begin{proposition}[拟凹是序数性质，凹是基数性质]\label{prop:qc-ordinal}
\textbf{(i)} 若 $u$ 拟凹，$f$ 严格递增，则 $f\circ u$ 拟凹；
\textbf{(ii)} 存在凹函数 $u$ 与严格递增 $f$ 使 $f\circ u$ 不凹。
\end{proposition}

\begin{proof}
\textbf{(i)} $f(u(\lambda y+(1-\lambda)z))\ge f(\min\{u(y),u(z)\})=\min\{f(u(y)),f(u(z))\}$，
第一步用 $u$ 拟凹加 $f$ 递增，第二步因 $\min$ 与严格递增映射可交换。

\textbf{(ii)} 取 $u(t)=\sqrt{t}$（$t\ge0$，凹）与 $f(s)=s^4$（$s\ge0$ 上严格递增）：
$f(u(t))=t^2$，二阶导数 $2>0$，凸而不凹。
\end{proof}

\begin{remark}[判别清单]\label{rem:ordinal-checklist}
判断一个效用性质是否“偏好可表达”，标准是：
该性质是否\textbf{在一切严格递增变换下保持}（命题~\ref{prop:ordinal} 的变换群）。
\begin{center}
\small
\begin{tabular}{lll}
\toprule
性质 & 序数/基数 & 偏好侧对应 \\
\midrule
递增 / 严格递增 & 序数 & 单调 / 强单调 \\
拟凹 / 严格拟凹 & 序数 & 凸 / 严格凸偏好 \\
一次齐次 + 递增变换 & 序数 & 位似性 \\
加和形式（存在某表示） & 序数 & 坐标独立性（§3.6） \\
线性计价物形式（存在某表示） & 序数 & 无财富效应（§3.7） \\
连续 & 序数 & 连续偏好 \\
凹 / 边际效用递减 & 基数 & 无偏好对应（表示依赖） \\
边际替代率 & 序数 & 无差异曲线斜率 \\
弹性类指标（如替代弹性） & 序数 & 无差异曲线曲率 \\
\bottomrule
\end{tabular}
\end{center}
“边际效用递减”是最常见的误用：$u=\sqrt{x}$ 与 $\tilde u=u^3=x^{3/2}$ 表示\textbf{同一偏好}，
前者边际效用递减、后者递增——序数理论中“边际效用”没有行为含义，
有意义的只是边际效用\textbf{之比}（MRS）。
\end{remark}

\subsection{回到词典序：翻译词典的存在性前提}

把 §3.1 的反例放回本节框架：词典序偏好对第一坐标强单调，
但不存在任何效用函数承载它的结构——\textbf{翻译词典生效的前提是表示存在}
（§3.1 的可数序稠密条件）。这正是本章把存在性（§3.1--§3.4）置于性质对应（§3.5--§3.7）之前的逻辑顺序。

% ============================================================
\section{从坐标独立性到可分性：可加表示定理}\label{sec:separability}
% ============================================================

本节回答：偏好满足什么\textbf{条件}，才允许把效用\textbf{拆开写}成
各坐标贡献之和？结论是充要条件：
\[
\boxed{\ \text{可加表示 }u(x,y)=u_1(x)+u_2(y)\ \Longleftrightarrow\ \text{坐标独立性}+\text{矩形条件（}+ \text{正则性）}\ }
\]
并且当因子数 $\ge 3$ 时，矩形条件可由坐标独立性自动推出。

\subsection{坐标独立性条件}

\begin{condition}[坐标独立性]\label{cond:indep}
设 $X=X_1\times X_2$，$\succeq$ 为 $X$ 上的弱序。称 $\succeq$ 满足\textbf{坐标独立性}，
当且仅当对任意 $x,x'\in X_1$ 与 $y,y'\in X_2$：
\[
(x,y)\succeq(x',y)\ \Longleftrightarrow\ (x,y')\succeq(x',y'),
\]
且对称地对第二坐标成立：
\[
(x,y)\succeq(x,y')\ \Longleftrightarrow\ (x',y)\succeq(x',y').
\]
即：固定第二坐标后对第一坐标的比较，与第二坐标固定在何处\textbf{无关}。
\end{condition}

\begin{remark}
坐标独立性允许我们把“$X_1$ 上的诱导偏好”良定义为：
$x\succeq_1 x'\iff$ 任取（从而所有）$y$，$(x,y)\succeq(x',y)$。
它是“可分”二字的全部序内容：\textbf{一个坐标内的排序不被另一个坐标的水平污染}。
\end{remark}

\begin{condition}[矩形条件（双消去条件）]\label{cond:thomsen}
对任意 $x_1,x_2,x_3\in X_1$ 与 $y_1,y_2,y_3\in X_2$：若
\[
(x_1,y_2)\sim(x_2,y_1)\quad\text{且}\quad (x_2,y_3)\sim(x_3,y_2),
\]
则 $(x_1,y_3)\sim(x_3,y_1)$。
\end{condition}

\begin{remark}[矩形条件的可加性含义]
若 $u=u_1+u_2$ 存在，前提两式给出
$u_1(x_1)+u_2(y_2)=u_1(x_2)+u_2(y_1)$ 与 $u_1(x_2)+u_2(y_3)=u_1(x_3)+u_2(y_2)$，
两式相加消去公共项得 $u_1(x_1)+u_2(y_3)=u_1(x_3)+u_2(y_1)$，即结论式。
故矩形条件是可加表示的\textbf{必要条件}——它正是“坐标贡献可以互相抵消”的序表达。
\end{remark}

\begin{example}[可分性的数值检验]
取 $X_1=X_2=\R_+$。
\textbf{(i) 可分}：$u(x,y)=\sqrt{x}+\sqrt{y}$。取三点 $x_1=1,x_2=4,x_3=9$ 与 $y_1=1,y_2=4,y_3=9$。
前提一 $(x_1,y_2)\sim(x_2,y_1)$：$u(1,4)=1+2=3=u(4,1)$，成立；
前提二 $(x_2,y_3)\sim(x_3,y_2)$：$u(4,9)=2+3=5=u(9,4)$，成立；
结论 $(x_1,y_3)\sim(x_3,y_1)$：$u(1,9)=1+3=4=u(9,1)$，成立——矩形闭合。
\textbf{(ii) 不可分}：$u(x,y)=xy$。前提一要求 $x_1y_2=x_2y_1$，取 $x_1=1,y_2=4,x_2=2,y_1=2$ 满足 $1\cdot4=2\cdot2$；
前提二要求 $x_2y_3=x_3y_2$，取 $y_3=8,x_3=1$ 得 $2\cdot8=1\cdot16$ 成立；
但结论要求 $x_1y_3=x_3y_1$ 即 $1\cdot8=1\cdot2$，不成立——矩形条件破裂，
正对应 $u$ 不能写成 $u_1(x)+u_2(y)$ 的加和形式（交叉偏导 $\partial^2u/\partial x\partial y=y+x\ne0$）。
\end{example}

\begin{center}
\begin{tikzpicture}[scale=1.1]
  % grid
  \foreach \x/\xl in {1/x_1, 3/x_2, 5/x_3}{
    \draw[MidBlue!30] (\x,0.6) -- (\x,4.6);
    \node at (\x,0.3) {$\xl$};
  }
  \foreach \y/\yl in {1/y_1, 2.6/y_2, 4.2/y_3}{
    \draw[MidBlue!30] (0.6,\y) -- (5.4,\y);
    \node at (0.3,\y) {$\yl$};
  }
  % premise 1: (x1,y2)~(x2,y1)
  \node[nodeblue] at (1,2.6) {}; \node[nodeblue] at (3,1) {};
  \draw[thick, AxisBlue] (1,2.6) -- (3,1) node[midway, above left] {$\sim$};
  % premise 2: (x2,y3)~(x3,y2)
  \node[nodeblue] at (3,4.2) {}; \node[nodeblue] at (5,2.6) {};
  \draw[thick, AxisBlue] (3,4.2) -- (5,2.6) node[midway, above left] {$\sim$};
  % conclusion: (x1,y3)~(x3,y1)
  \node[nodeblue] at (1,4.2) {}; \node[nodeblue] at (5,1) {};
  \draw[thick, dashed, AxisBlue] (1,4.2) -- (5,1) node[midway, above right] {$\sim$（结论）};
  \node[align=left] at (6.35,2.6) {两条实线是前提，\\ 虚线是推论：\\ 大矩形的对角\\ 也无差异};
\end{tikzpicture}
\end{center}
\figexplan{网格给出 $x_1,x_2,x_3$ 与 $y_1,y_2,y_3$ 构成的二维截面。两条实线是可加可分的两个前提：$(x_1,y_2)\sim(x_2,y_1)$ 与 $(x_2,y_3)\sim(x_3,y_2)$；虚线是由矩形条件推出的结论 $(x_1,y_3)\sim(x_3,y_1)$。若虚线不成立（矩形破裂），则不存在加和可分表示 $u_1(x)+u_2(y)$。}

\subsection{可加表示定理（两因子）}

\begin{definition}[本质坐标与可解性]\label{def:essential}
坐标 $X_1$ 称为\textbf{本质的}：存在 $x,x'\in X_1$ 与 $y$ 使 $(x,y)\succ(x',y)$
（该坐标并非摆设）。\textbf{可解性}：对任意 $x,x'\in X_1$（$x\succeq_1 x'$）与 $y,y'\in X_2$，
方程 $(x,y)\sim(x',z)$ 在 $X_2$ 中有解 $z$——坐标间的“价值兑换”总是可行的。
\end{definition}

\begin{theorem}[两因子可加表示定理]\label{thm:additive}
设 $X=X_1\times X_2$，$\succeq$ 为连续弱序，两坐标均本质，满足
坐标独立性（条件~\ref{cond:indep}）、矩形条件（条件~\ref{cond:thomsen}）与可解性，
且 $X_1,X_2$ 连通。则存在连续函数 $u_1:X_1\to\R$、$u_2:X_2\to\R$，使
\[
u(x,y)=u_1(x)+u_2(y)
\]
表示 $\succeq$；且 $\{u_1,u_2\}$ 在“共同正仿射变换”下唯一：
任何其他可加表示 $(\tilde u_1,\tilde u_2)$ 都满足
$\tilde u_i=a u_i+b_i$（$a>0$）。
\end{theorem}

\begin{proof}[证明（构造与关键步骤）]
\textbf{第 1 步：取基准点，定义坐标刻度并完成刻度校准。}
固定基准 $(x^*,y^*)\in X$。由定理~\ref{thm:continuous-rep}，
$\succeq$ 存在连续表示 $U$。定义
\[
u_1(x)\coloneqq U(x,y^*)-U(x^*,y^*),\qquad
u_2(y)\coloneqq U(x^*,y)-U(x^*,y^*).
\]
由坐标独立性，$(x,y^*)\succeq(x',y^*)\iff(x,y)\succeq(x',y)$ 对一切 $y$ 成立，
故 $u_1$ 表示诱导偏好 $\succeq_1$；同理 $u_2$ 表示 $\succeq_2$。

\textbf{刻度校准}（把 $X_1$ 刻度兑换为 $X_2$ 刻度）：
由可解性，对每个 $x\in X_1$ 取 $w(x)\in X_2$ 使
\[
(x,y^*)\ \sim\ \bigl(x^*,w(x)\bigr).
\]
由 $U$ 表示 $\succeq$ 与 $u_1,u_2$ 的定义，
\begin{equation}\label{eq:calib}
u_1(x)\ =\ U(x,y^*)-U(x^*,y^*)\ =\ U\bigl(x^*,w(x)\bigr)-U(x^*,y^*)\ =\ u_2\bigl(w(x)\bigr)-u_2(y^*).
\end{equation}
即：\textbf{第一坐标的价值差可以无差异地平移为第二坐标的价值差}。

\textbf{第 2 步：矩形条件给出基本导出关系。}
任取 $x,x'\in X_1$。对方程组
\[
(x,y^*)\sim\bigl(x^*,w(x)\bigr),\qquad
\bigl(x^*,w(x')\bigr)\sim(x',y^*)
\]
施用矩形条件~\ref{cond:thomsen}（取三个 $X_1$ 层 $x,x^*,x'$，
三个 $X_2$ 层 $w(x),y^*,w(x')$：前提恰为
$(x,y^*)\sim(x^*,w(x))$ 与 $(x^*,w(x'))\sim(x',y^*)$），得结论
\begin{equation}\label{eq:derived}
\bigl(x,w(x')\bigr)\ \sim\ \bigl(x',w(x)\bigr).
\end{equation}

\textbf{第 3 步：可加分解。}
考察映射 $\Phi(x,y)\coloneqq u_1(x)+u_2(y)$。由式~\eqref{eq:calib}，
$\Phi(x,y)=u_2(w(x))+u_2(y)-u_2(y^*)$。
由式~\eqref{eq:derived}，$(x,y)\sim(x',y')$ 当且仅当沿导出关系逐层平移后
两点的 $X_2$ 刻度之和相等：
\[
(x,y)\sim(x',y')
\ \Longleftrightarrow\
u_2\bigl(w(x)\bigr)+u_2(y)\ =\ u_2\bigl(w(x')\bigr)+u_2(y'),
\]
理由：比较 $(x,y)$ 与 $(x',y')$ 时，先用导出关系~\eqref{eq:derived}
把 $(x',y')$ 换为与之无差异的 $\bigl(x,w(x')\bigr)$ 型点，
使两比较点落在同一 $X_1$ 层 $x$ 上；同层比较只由第二坐标决定
（坐标独立性），归结为 $u_2$ 的比较；
而导出关系本身（再用一次矩形条件传递到任意 $y$ 层）表明
该比较恰由 $u_2(w(x))+u_2(y)$ 与 $u_2(w(x'))+u_2(y')$ 的大小决定。
代回式~\eqref{eq:calib} 即
\[
(x,y)\succeq(x',y')\ \Longleftrightarrow\ u_1(x)+u_2(y)\ \ge\ u_1(x')+u_2(y'),
\]
严格偏好对应严格不等号。连续性由 $U$ 连续与构造继承。

\textbf{第 4 步：唯一性。}
设 $\tilde u_1+\tilde u_2$ 是另一可加表示。
$\tilde u_1(\cdot)+\tilde u_2(y^*)$ 与 $u_1(\cdot)+u_2(y^*)$ 都表示 $\succeq_1$
且共用与 $X_2$ 的兑换率~\eqref{eq:calib}，故存在 $a>0,b_1$ 使 $\tilde u_1=a u_1+b_1$；
同理 $\tilde u_2=a' u_2+b_2$。两斜率必须相等：取 $x\succ_1 x'$ 与 $y'\succ_2 y$ 使
$u_1(x)+u_2(y)=u_1(x')+u_2(y')$（无差异对，存在性由可解性与连续性），
若 $a\ne a'$ 则 $\tilde u$ 在该对上严格不等，破坏无差异，矛盾。
\end{proof}

\subsection{推广：三因子以上矩形条件自动成立 \texorpdfstring{$\star$}{*}}

\begin{proposition}\label{prop:3factor}
设 $X=X_1\times X_2\times X_3$，$\succeq$ 为弱序，三个坐标均本质且满足
坐标独立性（任意两坐标固定时对第三坐标的比较不受影响）。
则矩形条件自动成立，从而（在可解性与连续性下）存在可加表示
$u(x_1,x_2,x_3)=u_1(x_1)+u_2(x_2)+u_3(x_3)$。
\end{proposition}

\begin{proof}[证明思路]
矩形条件的前提为 $(x_1,y_2)\sim(x_2,y_1)$ 与 $(x_2,y_3)\sim(x_3,y_2)$
（固定第三坐标为 $z_0$ 写成完整三元组）。
利用第三坐标的本质性取 $z\succ_3 z_0$，构造辅助点把两条无差异
“平移”到不同的第三坐标层，再用坐标独立性逐层搬回同一层，
传递性完成闭合。第三坐标提供了两因子情形缺失的\textbf{自由度}：
等值关系的“搬运”可以绕道第三维而不触碰矩形本身，
这正是 $n\ge3$ 时独立性蕴含矩形条件的机制。
\end{proof}

\subsection{坐标独立性的另一表述：逐行表示与递增包络}

坐标独立性还有一个与 §3.5 翻译词典直接衔接的等价表述。

\begin{proposition}[逐行表示定理]\label{prop:rowwise}
设 $\succeq$ 为 $X_1\times X_2$ 上的弱序，满足坐标独立性，
且 $X_1$ 上的诱导偏好可由 $v:X_1\to\R$ 表示、$v$ 的值域为区间。
则存在 $U:\R\times X_2\to\R$ 使
\[
u(x,y)=U(v(x),y)
\]
表示 $\succeq$，且 $U(\cdot,y)$ 对每个 $y$ \textbf{严格递增}。
\end{proposition}

\begin{proof}
\textbf{第 1 步（逐行构造）}：固定 $y$，$X_1\times\{y\}$ 上的偏好由坐标独立性等同于
$\succeq_1$，故 $u(x,y)\coloneqq U(v(x),y)$ 形式可行——
只要定义 $U(t,y)$ 为“$X_1\times\{y\}$ 行上、$v$ 值为 $t$ 的点的效用层级”。
由 1（坐标独立性：行内比较只依赖 $v$ 值）与 2（$v$ 值域为区间：每个 $t$ 都有原像），
\[
U(v(x),y)=u\bigl(v^{-1}(v(x)),y\bigr)=u(x,y).
\]

\textbf{第 2 步（$U(\cdot,y)$ 严格递增）}：
任取 $y\in X_2$ 与 $x,x'\in X_1$ 使 $v(x')>v(x)$。则
\[
\begin{aligned}
v(x')>v(x)
&\Longrightarrow (x',y_0)\succ(x,y_0) && (\text{$v$ 表示 }\succeq_1)\\
&\Longrightarrow (x',y)\succ(x,y) && (\text{坐标独立性})\\
&\Longrightarrow u(x',y)>u(x,y) && (\text{$u$ 表示 }\succeq)\\
&\Longrightarrow U(v(x'),y)>U(v(x),y) && (\text{第 1 步}).
\end{aligned}
\]
故 $U(\cdot,y)$ 对每个 $y$ 严格递增。
\end{proof}

\begin{remark}
命题~\ref{prop:rowwise} 的序内容：\textbf{第一坐标的“价值”可以用一个统一刻度 $v$ 度量，
第二坐标只决定这个刻度如何被“翻译”进总效用}（$U(\cdot,y)$ 是逐 $y$ 的翻译器）。
可分偏好正是“边际估值不串扰”的偏好：
效用 $u(x,y)=x+\sqrt{y}$ 中，$x$ 的边际效用恒为 $1$，与 $y$ 的水平无关；
反例 $u(x,y)=xy$ 中，$x$ 的边际效用为 $y$——坐标间互相“赋能”，坐标独立性失效：
$(2,1)\succ(1,1)$ 而比较 $(2,0)$ 与 $(1,0)$ 时二者无差异。
\end{remark}

% ============================================================
\section{理性偏好的结构性质：拟线性}\label{sec:quasilinear}
% ============================================================

拟线性是应用工作（产业组织、公共经济学、福利分析）中使用最频繁的效用形式。
本节给出它的\textbf{偏好侧完整刻画}：三个序条件 $\iff$ 线性计价物表示，
并推导其两大后果（无财富效应、福利的精确货币度量）。

\subsection{设定与三个计价物条件}

设 $X=\R\times Y$：第一坐标 $t\in\R$ 为\textbf{计价物}（数量可正可负、连续可分），
$Y$ 为一般集合（不必是欧氏空间），$\succeq$ 为 $X$ 上的弱序。

\begin{condition}[计价物有价值]\label{cond:num-valuable}
对一切 $y\in Y$ 与 $t,t'\in\R$：
\[
(t,y)\succeq(t',y)\ \Longleftrightarrow\ t\ge t'.
\]
同一 $y$ 下，计价物越多越好，且严格。
\end{condition}

\begin{condition}[补偿可行]\label{cond:compensable}
对任意 $y,y'\in Y$，存在 $t\in\R$ 使
\[
(0,y)\ \sim\ (t,y').
\]
任意两个 $Y$ 方案之间的价值差都能用计价物\textbf{恰好补偿}。
\end{condition}

\begin{condition}[无财富效应]\label{cond:no-wealth}
若 $(t,y)\succeq(t',y')$，则对一切 $d\in\R$：
\[
(t+d,\,y)\ \succeq\ (t'+d,\,y').
\]
给两个方案同时增加等量计价物，不改变比较结果——偏好对计价物的“初始存量”不敏感。
\end{condition}

\begin{example}[无财富效应的直觉]
若 $(10,\,y)\sim(0,\,y')$——“持有 10 单位计价物配 $y$”与“不持计价物配 $y'$”无差异，
则同时加 100：$(110,\,y)\sim(100,\,y')$ 仍无差异。
财富水平的变化不翻转、不破坏任何比较。
\end{example}

\subsection{拟线性表示定理}

\begin{theorem}[拟线性表示定理]\label{thm:ql}
设 $\succeq$ 为 $X=\R\times Y$ 上的弱序，满足
条件~\ref{cond:num-valuable}、\ref{cond:compensable}、\ref{cond:no-wealth}。
则存在 $v:Y\to\R$ 使
\[
u(t,y)\ =\ t+v(y)
\]
表示 $\succeq$。
反之，任何由 $u(t,y)=t+v(y)$ 形式表示的偏好都满足三个计价物条件。
\end{theorem}

\begin{proof}
\textbf{正方向。}
\textbf{第 1 步（定义 $v$）}：固定基准 $\bar y\in Y$。
由条件~\ref{cond:compensable}，对每个 $y\in Y$ 存在 $s\in\R$ 使
$(0,y)\sim(s,\bar y)$；此 $s$ \textbf{唯一}——
若 $(s,\bar y)\sim(0,y)\sim(s',\bar y)$，传递性给出 $(s,\bar y)\sim(s',\bar y)$，
由条件~\ref{cond:num-valuable} 得 $s=s'$。定义
\[
v(y)\ \coloneqq\ \text{唯一的}\ s\ \text{使}\ (0,y)\sim(s,\bar y).
\]

\textbf{第 2 步（把任意点折算到 $\bar y$ 行）}：
对任意 $(t,y)$，由 $(0,y)\sim(v(y),\bar y)$ 与条件~\ref{cond:no-wealth}
（取 $d=t$）得
\[
(t,\,y)\ \sim\ \bigl(t+v(y),\,\bar y\bigr).
\]

\textbf{第 3 步（表示性验证）}：
对任意 $(t,y),(t',y')\in X$，
\[
\begin{aligned}
(t,y)\succeq(t',y')
&\iff \bigl(t+v(y),\bar y\bigr)\succeq\bigl(t'+v(y'),\bar y\bigr)
   && (\text{第 2 步} + \text{传递性})\\
&\iff t+v(y)\ \ge\ t'+v(y')
   && (\text{条件~\ref{cond:num-valuable}}).
\end{aligned}
\]
即 $u(t,y)=t+v(y)$ 表示 $\succeq$。

\textbf{逆方向。}设 $u(t,y)=t+v(y)$ 表示 $\succeq$，逐条验证：
\textbf{条件~\ref{cond:num-valuable}}：$(t,y)\succeq(t',y)\iff t+v(y)\ge t'+v(y)\iff t\ge t'$。
\textbf{条件~\ref{cond:compensable}}：取 $t=v(y)-v(y')$，则
$u(0,y)=v(y)=t+v(y')=u(t,y')$，即 $(0,y)\sim(t,y')$。
\textbf{条件~\ref{cond:no-wealth}}：$(t,y)\succeq(t',y')\iff t+v(y)\ge t'+v(y')
\iff (t+d)+v(y)\ge(t'+d)+v(y')\iff(t+d,y)\succeq(t'+d,y')$。
\end{proof}

\begin{remark}[证明的结构]
三步证明做了一件事：\textbf{用条件~\ref{cond:compensable} 把每个 $y$ “定价”为 $v(y)$，
用条件~\ref{cond:no-wealth} 把定价沿计价物轴平移到任何财富水平，
用条件~\ref{cond:num-valuable} 在单一 $\bar y$ 行上完成所有比较}。
三条缺一，证明即断：无~\ref{cond:compensable} 则 $v$ 无定义；
无~\ref{cond:no-wealth} 则定价随财富漂移；
无~\ref{cond:num-valuable} 则 $\bar y$ 行上的比较无法数值化。
\end{remark}

\begin{proposition}[$v$ 的唯一性]\label{prop:ql-unique}
若 $t+v(y)$ 与 $t+\tilde v(y)$ 表示同一 $\succeq$，则 $\tilde v=v+\text{常数}$。
\end{proposition}

\begin{proof}
两者都表示时，$(0,y)\sim(0,y')\iff v(y)=v(y')\iff\tilde v(y)=\tilde v(y')$；
且 $(t,y)\sim(0,y')\iff t=v(y')-v(y)\iff t=\tilde v(y')-\tilde v(y)$，
故 $v(y')-v(y)=\tilde v(y')-\tilde v(y)$ 对一切 $y,y'$ 成立，
固定 $y$ 即得 $\tilde v-v$ 为常数。
\end{proof}

\begin{remark}[序数性辨析]
定理说“存在 $t+v(y)$ 形式的表示”，并\textbf{不}排除其他形式的表示：
$u(t,y)=t+v(y)$ 经严格递增变换 $f$ 后一般不再是拟线性形式，
却仍表示同一偏好（命题~\ref{prop:ordinal}）。
“拟线性”是偏好的\textbf{序性质}（存在某表示具此形式），
不是某个特定函数的\textbf{形式性质}。
\end{remark}

\subsection{两大后果}

\notehead{后果一：需求无财富效应（收入扩张路径平行于计价物轴）}

设 $Y=\R_+$（记其元素为 $q$），$v$ 可微、严格凹，价格结构为“每单位 $q$ 兑换 $p$ 单位计价物”，
预算约束 $t+pq=m$。效用最大化：
\[
\max_{q\ge0}\ \bigl[m-pq+v(q)\bigr]
\ \Longrightarrow\ v'(q^*)=p\quad(\text{内点解}).
\]
一阶条件\textbf{不含 $m$}：最优 $q^*$ 只由 $p$ 决定，收入 $m$ 的变化全部吸收到 $t$ 上。
收入扩张路径是一条竖直线（$q$ 恒为 $q^*$）——这就是“无财富效应”的需求侧表达。

\begin{center}
\begin{tikzpicture}[scale=1.05]
  \draw[axisline] (0,0) -- (6.2,0) node[below left, axislabel] {$q$};
  \draw[axisline] (0,0) -- (0,4.8) node[below left, axislabel] {$t$};
  % parallel ICs: t = c - sqrt(q)
  \draw[thick, MidBlue, domain=0.02:5.8, smooth, variable=\x] plot ({\x},{3.6-sqrt(\x)});
  \draw[thick, MidBlue, domain=0.02:5.8, smooth, variable=\x] plot ({\x},{2.6-sqrt(\x)});
  \draw[thick, MidBlue, domain=0.02:5.8, smooth, variable=\x] plot ({\x},{1.6-sqrt(\x)});
  \node at (5.0,1.6) {$u=c$};
  % vertical expansion path
  \draw[dashed, thick, AxisBlue] (2.6,0) -- (2.6,4.6) node[above, axislabel] {收入扩张路径};
  \node[nodeblue] at (2.6,3.6-1.612) {};
  \node[nodeblue] at (2.6,2.6-1.612) {};
  \node[nodeblue] at (2.6,1.6-1.612) {};
  \node[align=left] at (4.6,3.9) {无差异曲线互为\textbf{纵向平移}：\\ $t+v(q)=c\iff t=c-v(q)$，\\ 同一 $q$ 处斜率恒为 $v'(q)$};
\end{tikzpicture}
\end{center}
\figexplan{三条无差异曲线互为纵向平移：$t+v(q)=c$ 等价于 $t=c-v(q)$，故在同一 $q$ 处斜率恒为 $v'(q)$。竖直虚线是收入扩张路径：一阶条件 $v'(q^*)=p$ 不含收入 $m$，收入变化全部由计价物 $t$ 吸收，$q^*$ 不变——这就是无财富效应的几何表现。}

\begin{remark}[对照：正常品的收入扩张路径]
一般偏好（如 $u(q,t)=q^{1/2}t^{1/2}$）下，收入增加时 $q^*$ 与 $t^*$ 同时增加，
收入扩张路径弯曲向右上——存在财富效应。
拟线性是一般偏好的\textbf{局部近似}：当计价物支出占财富比重很大、
$q$ 的支出占比很小时，财富效应对 $q$ 的需求影响可忽略，
拟线性是“部分均衡分析”的效用基础。
\end{remark}

\notehead{后果二：福利的精确货币度量}

比较两个 $Y$ 方案 $y,y'$ 时，拟线性偏好给出\textbf{与财富无关}的补偿度量：
\[
(t,y)\succeq(t,y')\ \Longleftrightarrow\ v(y)\ge v(y'),
\]
且从 $y'$ 切换到 $y$ 的\textbf{补偿变化}恰为
\[
v(y)-v(y')\ =\ \text{使}\ (t,y')\sim\bigl(t-\bigl[v(y)-v(y')\bigr],\,y\bigr)\ \text{成立的计价物数量}.
\]
即：$v(y)-v(y')$ 是“为保持在 $y'$ 原状，愿意为换到 $y$ 支付的最大计价物数量”，
与初始财富 $t$ \textbf{无关}。一般偏好下补偿变化依赖财富水平，
$v$ 之差与财富无关正是拟线性被广泛用于福利计算（剩余度量）的原因。

\begin{example}[补偿变化的数值计算]\label{ex:ql-welfare}
设 $v(y)=\sqrt{y}$（$y$ 为某种公共品供给量），计价物 $t$ 为税收（可负，即补贴）。
比较两个方案：现状 $y'=4$（$v(y')=2$）与拟建项目 $y=9$（$v(y)=3$）。
由上述结论，从 $y'$ 切换到 $y$ 的补偿变化为
\[
v(y)-v(y')=3-2=1,
\]
即：为使纳税人在“维持现状 $y'$”与“接受项目 $y$”间无差异，需向其补偿 $1$ 单位计价物；
该数值与初始财富 $t$ 无关。若改用一般（非拟线性）偏好 $u(t,y)=t^{1/2}y^{1/2}$ 作比较，
对同一对 $(t,4)$ 与 $(t-1,9)$，无差异要求
$t^{1/2}\cdot2=(t-1)^{1/2}\cdot3$，解得 $t=9/5=1.8$ 时才无差异——
补偿量随 $t$ 变化，直观显示了财富效应如何破坏“精确货币度量”。
\end{example}

\subsection{拟线性与可分性的关系}

拟线性是可加可分（§3.6）的特例加限制：
$u(t,y)=t+v(y)$ 显然是可加形式 $u_1(t)+u_2(y)$，
额外要求 $u_1$ 为\textbf{恒等映射}（对计价物坐标“价值=数量”）。
在条件语言中：坐标独立性 + 矩形条件给出 $a\cdot u_1(t)+v(y)$ 型表示，
条件~\ref{cond:num-valuable} 与 \ref{cond:no-wealth} 进一步把 $u_1$ 钉死为线性
（无财富效应排除了 $u_1(t)=\sqrt{t}$ 之类的弯曲刻度）。
因此拟线性偏好自动满足坐标独立性与矩形条件：
$(t,y)\succeq(t',y)\iff t\ge t'$ 与 $(t,y')\succeq(t',y')$ 同真同假——
第一坐标的比较与 $y$ 固定在何处无关。

\subsection{可微情形 \texorpdfstring{$\star$}{*}}

设 $Y=\R^{n-1}_+$，$v$ 二阶连续可微。$u(t,y)=t+v(y)$ 的梯度为 $(1,\nabla v(y))$，
边际替代率
\[
\mathrm{MRS}_{t,i}(t,y)=\frac{\partial u/\partial y_i}{\partial u/\partial t}=\frac{\partial v}{\partial y_i}(y)
\]
\textbf{与 $t$ 无关}——无财富效应的可微表达。
拟凹性等价于 $v$ 凹：
$u$ 的海塞矩阵为 $\begin{pmatrix}0&0\\0&D^2v\end{pmatrix}$，
其加边海塞行列式条件归结为 $D^2v$ 负半定。
凸偏好（命题~\ref{prop:convex-qc}）$\iff$ $v$ 凹，
严格凸 $\iff$ $v$ 严格凹——拟线性把偏好凸性完全压缩为 $v$ 的凹性。

% ============================================================
\section{行为批判：公理体系的系统性失效}\label{sec:behavioral}
% ============================================================

理性选择框架（完备 + 传递 + 效用最大化）是分析基准，
但大量受控实验与现场数据表明：人的实际选择\textbf{系统性地}偏离公理。
批判分两类：
\textbf{情境依赖}（偏好随选择集、表述方式、社会情境改变——偏好并非预先存在）
与\textbf{非深思选择}（选择不经由两两比较与最大化，而由直觉、启发式驱动）。
每条批判都精确对应某条公理或某条结构性质的失效。

\subsection{批判一：情境依赖——诱饵效应}

\begin{example}[诱饵效应]\label{ex:decoy}
两属性商品空间：质量（纵轴）与价格（横轴，左好右差）。
\begin{itemize}
  \item 选项 $A$：高质量、高价格；
  \item 选项 $B$：低质量、低价格；
  \item 诱饵 $C$：质量略低于 $A$、价格高于 $A$（被 $A$ 严格支配）。
\end{itemize}
选择集 $\{A,B\}$ 中，$A$ 与 $B$ 各获约一半选择；
加入\textbf{无人选择的}诱饵 $C$ 后，选 $A$ 的比例显著上升。
\end{example}

\begin{center}
\begin{tikzpicture}[scale=1.1]
  \draw[axisline] (0,0) -- (6.4,0) node[below left, axislabel] {价格（向右递增）};
  \draw[axisline] (0,0) -- (0,4.6) node[below left, axislabel] {质量};
  \node[nodeblue, label=above left:{$A$}] at (3.4,3.4) {};
  \node[nodeblue, label=above left:{$B$}] at (1.2,1.1) {};
  \node[nodeblue, label=below right:{$C$}] at (4.6,3.0) {};
  % dominance region of A over C
  \draw[dashed, MidBlue] (3.4,3.4) -- (4.6,3.4) -- (4.6,3.0);
  \node[align=left] at (3.35,4.25) {$A$：高质高价，$B$：低质低价，互不占优};
  \node[align=left] at (4.45,1.9) {$C$（诱饵）：质量低于 $A$、\\ 价格高于 $A$，被 $A$ 严格支配};
\end{tikzpicture}
\end{center}
\figexplan{横轴价格向右递增、纵轴质量向上递增。$A$ 高质高价、$B$ 低质低价，二者互不占优；诱饵 $C$ 质量低于 $A$ 且价格高于 $A$，因而被 $A$ 严格支配（虚线折线标出支配关系）。$C$ 本身无人选择，却显著提高 $A$ 的选择份额，直接违反理性框架的正则性。}

\begin{proposition}[诱饵效应违反正则性]\label{prop:decoy}
记 $P(x,S)$ 为从选择集 $S$ 中选中 $x$ 的（群体）频率。
理性框架的\textbf{正则性}要求：$T\subseteq S\Rightarrow P(x,T)\ge P(x,S)$——
选择集扩大不会提升原有元素被选中的份额。
诱饵实验中 $P\bigl(A,\{A,B,C\}\bigr)>P\bigl(A,\{A,B\}\bigr)$，
直接违反正则性；在确定性选择函数语言下，
这对应于违反显示偏好理论中的\textbf{无关选项独立性}（IIA）。
\end{proposition}

\begin{remark}[机制]
诱饵提供了\textbf{局部比较基准}：$A$ 对 $C$ 的明显优势被归因为 $A$ 的“品质”，
而 $A$ 与 $B$ 之间需要跨属性权衡（困难比较）。
选择依赖“哪个比较容易论证”，而非稳定的内在排序——
偏好在选择时刻被\textbf{构造}，而非被\textbf{揭示}。
\end{remark}

\subsection{批判二：框架效应}

\begin{example}[生存率与死亡率]\label{ex:framing}
同一项治疗方案的两种表述：
\begin{itemize}
  \item 表述一：“治疗后存活率为 $90\%$”；
  \item 表述二：“治疗后死亡率为 $10\%$”。
\end{itemize}
两种表述在信息上\textbf{完全等价}，但表述一下接受率显著更高。
类似地，“含脂 $25\%$”与“瘦肉 $75\%$”的同一食品获得不同评价。
\end{example}

\begin{remark}[违反的公理：描述不变性]
理性框架隐含\textbf{描述不变性}：选择只依赖选项的结果结构，
与表述方式无关。框架效应表明描述本身进入偏好——
同一结果集在不同“参照点 + 损益编码”下被评估为不同方案。
这直接动摇效用表示的前提：$u$ 的定义域不再是“客观结果”，
而是“被框架化后的结果”，$\succeq$ 在客观结果集上甚至\textbf{无法良定义}。
\end{remark}

\subsection{批判三：偏好反转}

\begin{example}[定价与选择的矛盾]\label{ex:reversal}
两个不确定性选项：$P$ 型（高概率、小收益）与 $\$$ 型（低概率、大收益）。
受试者在\textbf{直接选择}中选 $P$ 型，
但对两者的\textbf{定价}（最低愿售价）却给 $\$$ 型定更高价。
若偏好稳定，$P\succ\$$ 意味着 $P$ 的定价应更高——观察到的是系统性反转。
\end{example}

\begin{remark}[违反的公理：程序不变性]
选择与定价是同一潜在偏好的两种\textbf{导出方式}，程序不变性要求它们一致。
反转表明：不同引出程序激活不同的认知过程
（选择任务突出概率维度、定价任务突出金额维度），
产出不同的“偏好”——不存在一个同时解释两者的稳定 $\succeq$。
\end{remark}

\subsection{批判四：非深思选择——启发式、默认与现状偏差}

\begin{example}[启发式选择]
购买复杂商品时，以“选最知名品牌”“选销量最高”“选中间价位”
代替逐属性比较——这些规则\textbf{不遍历}全部选项，
不满足“对所有备选作两两比较”的完备性操作。
\end{example}

\begin{example}[默认效应与现状偏差]\label{ex:default}
同一保险计划，设为\textbf{默认加入}（需主动退出）时参保率，
远高于设为\textbf{默认退出}（需主动加入）时。
若偏好稳定，默认设置（一个不改变任何选项可行性的“框”）不应改变选择。
现状偏差的形式化：存在\textbf{参照依赖}——
以现状 $r$ 为参照点时 $\succeq_r$，以 $r'$ 为参照点时 $\succeq_{r'}$，
二者不必一致；特别地，“失去”的权重大于等量“得到”（损失厌恶），
使无差异曲线在参照点处出现\textbf{折点}。
\end{example}

\begin{center}
\begin{tikzpicture}[scale=1.05]
  \draw[axisline] (0,0) -- (6.2,0) node[below left, axislabel] {商品 1};
  \draw[axisline] (0,0) -- (0,4.4) node[below left, axislabel] {商品 2};
  \node[nodeblue, label=below right:{参照点 $r$}] at (3.0,2.0) {};
  % kinked IC through r
  % 折点无差异曲线：左支陡（损失域）、右支平（收益域），两支均向右下倾斜
  \draw[thick, MidBlue] (0.4,3.72) -- (3.0,2.0) -- (5.9,1.18);
  \node[figlabel] at (1.05,3.98) {无差异曲线};
  % 虚线：无损失厌恶时的直线基准（过 r，斜率与损失域相同）
  \draw[dashed, MidBlue] (0.4,2.74) -- (5.5,0.37);
  \node[align=left,figlabel] at (1.30,1.80) {损失域（左上）：斜率陡峭，\\放弃一单位需更大补偿};
  \node[align=left,figlabel] at (4.70,2.60) {收益域（右下）：斜率平缓，\\同等收益只需较小付出};
\end{tikzpicture}
\end{center}
\figexplan{无差异曲线在参照点 $r$ 处折曲，且两支均向右下倾斜：左支（损失域，失去商品 1、获得商品 2）斜率陡，右支（收益域）斜率平缓。虚线是“若无损失厌恶”时的直线基准，其斜率与损失域相同；实际曲线在 $r$ 左侧高于该基准（更陡）、在右侧也高于该基准（更平）。两支斜率之比即损失厌恶程度：放弃一单位所需的补偿大于获得同量带来的收益。}

\begin{remark}[禀赋效应：买卖价差]
同一物品，拥有者的最低愿售价（WTA）系统性高于非拥有者的最高愿购价（WTP）。
在 §3.7 拟线性框架下，补偿变化与等值变化应\textbf{相等}（无财富效应时）；
即使允许财富效应，观察到的价差也远超收入效应可解释的幅度。
参照依赖（$r$=“已拥有”对 $r'$=“未拥有”）给出统一解释：
“卖出”被编码为损失，“买入”被编码为得到。
\end{remark}

\subsection{批判五：锚定与偏好构建}

\begin{example}[锚定]
先让受试者写下证件号码后两位，再问其愿意为某商品支付多少：
号码大者出价系统性更高——一个\textbf{明显无关}的数字影响了估值。
\end{example}

\begin{remark}
锚定是“偏好构建论”的极端证据：估值过程从最近激活的数字出发做不充分调整，
说明许多“偏好”没有预先存储的数值，而是在被询问的当下\textbf{临时构造}。
这与理性框架的基本设定（偏好是先于选择的稳定序）正面冲突。
\end{remark}

\subsection{效应与公理的对照}

\begin{center}
\small
\begin{tabular}{lll}
\toprule
行为效应 & 被违反的公理 / 性质 & 失效机制 \\
\midrule
诱饵效应 & 正则性 / IIA & 局部比较基准改变相对吸引力 \\
折中效应 & 正则性 / IIA & 选项因“居中”而获益 \\
框架效应 & 描述不变性 & 损益编码随表述改变 \\
偏好反转 & 程序不变性 & 不同引出程序激活不同维度权重 \\
默认 / 现状偏差 & 偏好稳定性 & 参照点进入偏好本身 \\
禀赋效应（WTA$>$WTP） & 补偿对称性（§3.7） & 损失权重 $>$ 得到权重 \\
锚定 & 偏好预先存在性 & 估值当下构造、不充分调整 \\
冲动选择 & 深思比较假设 & 启发式 / 直觉替代权衡 \\
\bottomrule
\end{tabular}
\end{center}

\begin{remark}[文献线索]
本章各效应的奠基性实验与文献：诱饵（吸引）效应源于 Huber, Payne \& Puto (1982) 的“占优违背”实验，并经 Tversky \& Simonson (1993) 系统化为 attraction effect；框架效应由 Tversky \& Kahneman (1981) 在“亚洲疾病问题”中确立；偏好反转由 Lichtenstein \& Slovic (1971) 发现、Grether \& Plott (1979) 以受控实验严格复核；默认 / 现状偏差的经典现场证据见 Johnson \& Goldstein (2003) 对欧洲器官捐献制度的跨国比较，Thaler \& Benartzi (2004) 的“明天储蓄更多”（Save More Tomorrow）计划则是其政策应用；禀赋效应由 Thaler (1980) 提出、Kahneman, Knetsch \& Thaler (1990) 以“咖啡杯实验”给出干净识别；锚定效应见 Tversky \& Kahneman (1974)，Ariely, Loewenstein \& Prelec (2003) 用“随机数字轮盘”排除了“理性锚定”的替代解释；上述效应的统一理论框架是 Kahneman \& Tversky (1979) 的\textbf{前景理论}（参照依赖 + 损失厌恶 + 概率权重）。
\end{remark}

\subsection{综合：理性基准与行为视角}

\begin{center}
\small
\begin{tabular}{lll}
\toprule
 & 理性选择模型 & 行为视角 \\
\midrule
偏好 & 稳定、预先存在 & 选择时刻被构造 \\
决策规则 & 两两比较 + 效用最大化 & 启发式、直觉、局部比较 \\
对框架的敏感性 & 无（描述不变） & 系统性敏感 \\
参照点 & 无 & 核心（损失厌恶） \\
财富效应 & 拟线性情形为零 & 普遍存在且随情境变化 \\
理论角色 & 基准与近似 & 对偏离方向的解释与预测 \\
\bottomrule
\end{tabular}
\end{center}

\begin{remark}[如何对待批判]
批判并未使理性框架失效，而是标定了它的\textbf{适用边界}：
(i) 作为\textbf{基准}——偏离是有方向的（诱饵使目标选项份额上升），
理性模型给出“无偏离”的零点，行为模型给出偏离项；
(ii) 作为\textbf{近似}—— stakes 高、学习机会多、激励强的环境中，
选择更接近公理；一次性、低激励、高复杂度的选择中偏离最大；
(iii) 公理的可检验化（第二章的显示偏好路线）使“偏离”可以被度量：
对每个选择数据集，可计算其违反 SARP 的程度（不一致指数），
把“理性程度”从是非判断变为连续变量。
\end{remark}

\begin{remark}[本章收束]
从公理到表示（§3.1--§3.4），从表示到结构（§3.5--§3.7），再到结构的边界（§3.8）：
效用函数是\textbf{序信息的压缩编码}，其存在性由可数序稠密条件精确刻画，
其特殊形式（可加、拟线性）由相应的序条件精确刻画，
而其失效处（情境、框架、参照点）同样可以在公理语言中精确指认。
这套“条件 $\leftrightarrow$ 表示”的方法论，是后续消费者理论、
不确定情形下效用理论（期望效用及其批判）的直接模板。
\end{remark}

\clearpage
\section{拓展：表示定理的证明策略与可计算含义}

\subsection{从序关系到数值编码}
效用表示的核心并不是给每个方案随意贴一个数字，而是寻找保持序关系的映射。若 $x\succ y$，必须有 $u(x)>u(y)$；若 $x\sim y$，必须有 $u(x)=u(y)$。因此，构造过程首先处理无差异类，再对商集排序。

\begin{definition}[商集与严格序]
令 $I$ 为无差异关系，定义 $X/I$ 为无差异类的集合。若偏好完备且传递，则在适当条件下可在 $X/I$ 上定义严格序：$[x]>[y]$ 当且仅当 $x\succ y$。
\end{definition}

\begin{figure}[H]\centering
\begin{tikzpicture}[font=\scriptsize,>=stealth]
\node[draw,rounded corners,minimum width=2.4cm,minimum height=1cm,fill=gray!10] (x) at (0,0) {$X$：方案集合};
\node[draw,rounded corners,minimum width=2.6cm,minimum height=1cm,fill=gray!20] (q) at (4,0) {$X/I$：无差异类};
\node[draw,rounded corners,minimum width=2.6cm,minimum height=1cm,fill=gray!30] (r) at (8,0) {$u(X)\subseteq\mathbb R$};
\draw[->,thick] (x)--(q) node[midway,above] {合并 $x\sim y$};
\draw[->,thick] (q)--(r) node[midway,above] {保持严格序};
\node[align=center] at (4,-1.5) {先处理关系结构，再进行数值编码\\数字本身没有基数含义};
\end{tikzpicture}
\caption{效用表示的两步构造}
\end{figure}
\figexplan{左侧是原始方案集合，中间把无差异方案压缩为等价类，右侧才是实数值编码。箭头说明效用表示首先是序结构的嵌入，而不是对满意度的物理测量。}

\subsection{单调变换与识别边界}
若 $u$ 表示偏好，且 $\phi$ 严格递增，则 $\phi\circ u$ 表示同一偏好。证明直接来自
$u(x)\ge u(y)\Leftrightarrow\phi(u(x))\ge\phi(u(y))$。因此，实证研究中由选择数据通常只能识别效用的序数类，而不能识别零点、单位或曲率。

\begin{proposition}[严格递增变换不改变最优集合]
对任意非空可行集 $B$，若 $\phi$ 严格递增，则
\[
\arg\max_{x\in B}u(x)=\arg\max_{x\in B}\phi(u(x)).
\]
\end{proposition}
\begin{proof}
若 $x^*$ 最大化 $u$，则对任意 $x\in B$ 有 $u(x^*)\ge u(x)$，严格递增性给出 $\phi(u(x^*))\ge\phi(u(x))$。反向同理。
\end{proof}

\subsection{连续性、闭图与极限偏好}
令 $x^n\to x$ 且 $y^n\to y$。若每期 $x^n\succeq y^n$，偏好关系的闭性要求极限仍满足 $x\succeq y$。闭性是把有限样本推向连续状态空间的桥梁，也是连续效用表示存在的重要技术条件。

\begin{figure}[H]\centering
\begin{tikzpicture}[x=1cm,y=0.75cm,font=\scriptsize]
\draw[->] (0,0)--(7,0) node[right] {$n$};\draw[->] (0,0)--(0,4) node[above] {比较结果};
\draw[dashed] (0,2)--(7,2) node[right] {极限边界};
\draw[thick] (0.5,3.1)--(1.5,2.8)--(2.5,2.6)--(3.5,2.35)--(4.5,2.15)--(5.5,2.02)--(6.5,2.0);
\node[fill=white] at (3.2,3.45) {$x^n\succeq y^n$ 对所有 $n$};
\end{tikzpicture}
\caption{闭性把有限比较延伸到极限}
\end{figure}
\figexplan{曲线表示一列比较强度逐渐靠近边界。若偏好关系闭，则所有有限期成立的弱偏好在极限点仍成立；若不闭，极限可能落在关系图之外。}

\subsection{拟凹性与凸偏好的几何对应}
效用函数拟凹当且仅当其上水平集 $\{x:u(x)\ge c\}$ 为凸集。经济含义是：若两个消费束都至少达到 $c$，它们的混合束也至少达到 $c$，体现对多样化的偏好。

\begin{theorem}[拟凹性的等价表述]
对连续函数 $u$，以下命题等价：
\begin{enumerate}
\item $u$ 拟凹；
\item 对任意 $c$，上水平集 $\{x:u(x)\ge c\}$ 凸；
\item 对任意 $x,y$ 和 $t\in[0,1]$，有 $u(tx+(1-t)y)\ge\min\{u(x),u(y)\}$。
\end{enumerate}
\end{theorem}
\begin{proof}
(1) 与 (3) 是定义的两种写法。(3) 推出 (2)：若 $u(x),u(y)\ge c$，则右侧至少为 $c$，所以混合束仍在上水平集内。(2) 推出 (3)：令 $c=\min\{u(x),u(y)\}$，凸性即得。
\end{proof}

\subsection{计算练习：离散权数构造}
设 $X=\{a,b,c,d\}$，偏好为 $a\succ b\sim c\succ d$。令
$u(a)=3,u(b)=u(c)=2,u(d)=0$。任何严格递增变换均给出同一排序，例如 $v=2u+5$，但 $v(a)-v(b)=2$ 并没有“偏好强度”解释。该例说明序数效用中的差值和比值不可比较。

\subsection{本章小结}
存在性定理回答“能否表示”，构造定理回答“如何表示”，性质对应回答“表示后能否保留经济结构”。三者必须分开：一个函数可能表示排序，却不保留凸性、可加性或拟线性结构；附加结构必须由附加公理或可检验限制提供。

\subsection{表示定理的例子与证明核对}

\begin{example}[有限序的显式表示]
设 $X=\{a,b,c,d,e\}$ 且 $a\succ b\sim c\succ d\succ e$。把无差异类分为
$[a]=\{a\}$、$[b]=\{b,c\}$、$[d]=\{d\}$、$[e]=\{e\}$，依次赋值 $3,2,1,0$，则
$u(a)=3$、$u(b)=u(c)=2$、$u(d)=1$、$u(e)=0$。验证任意一对方案即可看到：严格顺序对应严格数值，不同类之间没有重复值，同一类内部数值相同。
\end{example}

\begin{proposition}[单调变换保持表示]
若 $u$ 表示 $\succeq$，$f$ 严格递增，则 $f\circ u$ 也表示 $\succeq$。
\end{proposition}
\begin{proof}
严格递增意味着对任意 $a,b$，$a\ge b$ 当且仅当 $f(a)\ge f(b)$。因此
$x\succeq y\Leftrightarrow u(x)\ge u(y)\Leftrightarrow f(u(x))\ge f(u(y))$，恰为表示定义。
\end{proof}

\begin{example}[同一偏好的三种写法]
对 $u(x_1,x_2)=x_1+x_2$，函数 $v=2u+1$ 与 $w=\ln(1+u)$ 在定义域上都严格递增变换，因此三者产生相同的无差异集合与最优选择。不能据此比较 $u(x)-u(y)$ 与 $v(x)-v(y)$ 的大小。
\end{example}

\subsection{可数构造的细节}
对可数序稠密集 $D=\{d_1,d_2,\ldots\}$，逐个插入有理数。插入 $d_n$ 时，若已有点在其下方和上方，取两者有理数坐标的中点；若只有一侧有界，则向该侧取足够远的有理数。每一步只改变新点坐标，不改变已有点之间的顺序。由归纳法，最终映射保持严格序。

\begin{proposition}[中点构造保持序]
若 $a<b$ 且 $a,b\in\Q$，则 $(a+b)/2\in\Q$ 且 $a<(a+b)/2<b$。
\end{proposition}
\begin{proof}
有理数对加法和除以非零整数封闭，故中点为有理数；两边不等式由 $b-a>0$ 直接得到。
\end{proof}

\begin{example}[三次插入]
先把 $d_1$ 映射为 $0$，把 $d_2\succ d_1$ 映射为 $1$。若第三个元素满足 $d_2\succ d_3\succ d_1$，则取 $g(d_3)=1/2$。第四个元素若位于 $d_3$ 与 $d_2$ 之间，则取 $3/4$。每次取中点，都保持已有坐标的相对顺序。
\end{example}

\subsection{可加表示的检验}
若偏好具有坐标独立性与矩形条件，可写为 $u(x_1,x_2)=v_1(x_1)+v_2(x_2)$。检验时先固定第二坐标比较第一坐标，再固定第一坐标比较第二坐标，最后用矩形条件排除交叉项。该步骤解释了为什么“每一维都可单独排序”仍不足以推出可加性：还需要不同坐标变化之间的替代关系一致。

\begin{remark}
表示定理的证明不仅回答存在性，也揭示限制来源：有限集依赖排序，可数集依赖稠密插入，连续消费空间依赖拓扑连续性，可加表示还依赖跨坐标一致性。不同结论不能互相替代。
\end{remark}

\subsection{更多例子：连续性、凸性与拟线性}

\begin{example}[线性效用的表示不唯一]
若 $u(x)=x_1+x_2$，则 $v(x)=\exp(x_1+x_2)$ 表示同一偏好，但 $v$ 的边际效用随消费水平同时增长。由此可见，边际效用大小不是序数表示的内在属性；只有无差异集合和选择集合在严格递增变换下保持不变。
\end{example}

\begin{proposition}[严格单调变换保持拟凹性]
若 $u$ 拟凹且 $f$ 严格递增，则 $f\circ u$ 拟凹。
\end{proposition}
\begin{proof}
对任意 $a$，有
$\{x:f(u(x))\ge a\}=\{x:u(x)\ge f^{-1}(a)\}$，其中只需在 $f$ 的值域内理解逆函数。右侧是 $u$ 的上水平集，为凸集，故左侧也凸。
\end{proof}

\begin{example}[凸偏好与非凸偏好]
$u(x)=\sqrt{x_1}+\sqrt{x_2}$ 的上水平集为凸集，体现多样化偏好；$u(x)=\max\{x_1,x_2\}$ 的上水平集为两个半空间的并，一般不凸，体现专业化偏好。两者都可表示，但表示存在不意味着偏好具有凸性。
\end{example}

\begin{example}[拟线性效用]
令 $u(m,z)=m+v(z)$，其中 $m$ 是计价物。若 $v$ 可微，则对 $z$ 的边际支付意愿为 $v'(z)$，不依赖 $m$。这正是拟线性结构的政策含义：收入变化不改变非货币品的边际评价，直到角点或收入效应重新进入模型。
\end{example}

\begin{proposition}[对角线编号的唯一性]
若偏好连续且强单调，则每个无差异类与对角线 $\{t\mathbf 1:t\ge0\}$ 至多交一次。
\end{proposition}
\begin{proof}
若 $t_2>t_1$ 且 $x\sim t_1\mathbf1\sim t_2\mathbf1$，则传递性给出 $t_1\mathbf1\sim t_2\mathbf1$。但强单调性给出 $t_2\mathbf1\succ t_1\mathbf1$，矛盾。因此交点至多一个。连续性与边界符号变化共同保证交点存在。
\end{proof}

\subsection{例题：从效用函数反推偏好性质}
给定 $u=x_1^\alpha x_2^{1-\alpha}$，$0<\alpha<1$。一阶偏导均为正，故严格单调；对数变换后 $\ln u=\alpha\ln x_1+(1-\alpha)\ln x_2$，其 Hessian 在正象限上负半定意义下给出拟凹性。于是无差异曲线向右下倾斜且向原点凸。若把指数改为负权重，则单调性和凸性均需重新检查，不能从“乘法形式”直接套结论。

\subsection{例题：从表示到选择预测}

\begin{example}[CES 偏好的替代弹性]
令 $u(x)=\left[\alpha x_1^\rho+(1-\alpha)x_2^\rho\right]^{1/\rho}$。对任意给定效用水平，边际替代率为
$\mathrm{MRS}=\frac{\alpha}{1-\alpha}(x_1/x_2)^{\rho-1}$。当 $\rho<1$ 时，随 $x_1/x_2$ 上升，MRS 下降，体现凸偏好；参数 $\rho$ 只在给定函数形式下有解释，严格递增变换不会改变选择，但会改变边际效用数值。
\end{example}

\begin{proposition}[位似偏好的需求比例]
若效用函数一次齐次，价格与收入同时乘以 $t>0$，最优消费束不变；若只收入乘以 $t$，需求按比例扩大。
\end{proposition}
\begin{proof}
一次齐次性给出 $u(tx)=tu(x)$，因此同一条射线上的束具有按比例变化的效用。预算约束在价格与收入同时乘以 $t$ 时化为原约束乘以同一因子，约束集合不变，故选择不变。只收入扩大时，若 $x^*$ 可行，则 $tx^*$ 处于新预算集；齐次性与严格拟凹性保证最优比例保持。
\end{proof}

\begin{example}[效用表示与福利判断的区别]
两个消费者的选择都可由 $u_i(x)=x$ 表示，但把消费者甲效用提高 1 与把消费者乙效用提高 1 是否同等重要，不由表示定理决定。社会福利函数需要额外给出权重或分配原则；个体效用的序数性不能自动生成社会比较。
\end{example}

\begin{example}[连续性检验]
若 $u(x)=x_1+x_2$，上水平集 $\{x:u(x)\ge c\}$ 和下水平集 $\{x:u(x)\le c\}$ 都是闭半空间，故偏好连续。若定义 $u(x)=1$ 当 $x_1>0$、$u(x)=0$ 当 $x_1=0$，则在边界处发生跳跃，连续性失败；连续效用表示定理不能直接应用。
\end{example}

\begin{remark}
本章的例题应按“偏好公理—效用表示—函数性质—选择预测”的顺序阅读。先判断函数能否表示，再判断它是否单调、拟凹或位似，最后才讨论需求和福利，避免把解析形式直接当成经济含义。
\end{remark}

\subsection{同一例题的四步核验：表示、性质、需求与比较静态}
考虑两商品 CES 形式
\[
 u(x_1,x_2)=\left[\alpha x_1^\rho+(1-\alpha)x_2^\rho\right]^{1/\rho},
 \qquad 0<\alpha<1,\quad \rho<1,\ \rho\ne0.
\]
第一步只讨论表示：由于 CES 函数连续，且在 $x_i>0$ 时严格递增，
其上水平集和下水平集均可用极限序列检验闭性，因此它确实承载一个连续偏好。
这一结论并不依赖具体价格，也不等同于“需求已经被求出”。

第二步核验凸性。对固定的效用水平 $\bar u$，将约束写成
\[
 \alpha x_1^\rho+(1-\alpha)x_2^\rho=\bar u^\rho.
\]
隐式求导得到
\[
 \frac{dx_2}{dx_1}
 =-\frac{\alpha}{1-\alpha}
   \left(\frac{x_1}{x_2}\right)^{\rho-1}.
\]
当 $\rho<1$ 时，沿曲线向右移动，绝对斜率下降；因此曲线向原点凸，
对应的上水平集是凸集。这里必须注意：斜率下降是局部微分判断，
凸偏好还需要由整个上水平集的集合性质确认，不能只凭一段曲线外观下结论。

第三步求需求。拉格朗日函数为
\[
 \mathcal L=\left[\alpha x_1^\rho+(1-\alpha)x_2^\rho\right]^{1/\rho}
 +\lambda(m-p_1x_1-p_2x_2).
\]
一阶条件相除后得到
\[
 \frac{\alpha}{1-\alpha}\left(\frac{x_1}{x_2}\right)^{\rho-1}
 =\frac{p_1}{p_2},
\]
从而
\[
 \frac{x_1}{x_2}
 =\left[\frac{\alpha p_2}{(1-\alpha)p_1}\right]^{1/(1-\rho)}.
\]
将该比例与预算方程联立，即可得到两项 Marshallian 需求。这个推导显示，
价格比决定的是消费比例，收入决定的是共同尺度；后者正是位似性的需求含义。

第四步作反事实核验。若 $p_1$ 上升而 $p_2,m$ 不变，则上式中的 $x_1/x_2$
下降；若 $m$ 与两种价格同时乘以 $t>0$，比例不变且预算约束整体缩放，
需求束保持不变。由此可把“函数形式—偏好性质—选择预测—比较静态”四层
严格区分：前一层的结论不能未经证明直接替代后一层。

\begin{example}[线性与 CES 的边界比较]
取 $\rho\to1$，CES 的替代性增强，等值线趋向直线，最优解可能转为角点或整段；
取 $\rho\to0$，经适当归一化后趋向 Cobb--Douglas 形式，需求比例由参数和价格共同决定；
取 $\rho\to-\infty$，等值线趋向完全互补的折线，最优点由两个固定比例约束共同确定。
这些是同一函数族内部的边界比较，用来核对图形、符号和极限，不是另行引入新的理论主题。
\end{example}

\subsection{例题核验表：不要跳过中间的逻辑层}
对任意给定的效用表达式，建议按以下顺序核对。首先检查定义域：若函数含有分数次幂、对数或负指数，必须先说明 $x_i>0$ 还是允许边界点。其次检查表示关系：需要验证 $x\succeq y$ 当且仅当 $u(x)\ge u(y)$，而不是只验证几个数值例子。再次检查性质：单调性由偏导数或集合包含关系确认，凸偏好由上水平集确认，位似性则由无差异曲线沿射线的比例结构确认。最后才把价格和收入代入最优化问题。

以 Cobb--Douglas 形式 $u=x_1^\alpha x_2^{1-\alpha}$ 为例，定义域边界 $x_i=0$ 不能被简单略去。若 $0<\alpha<1$，内点偏导数存在且为正，但在边界处某一偏导可能发散；这不影响连续性，却会影响直接使用内点一阶条件的合法性。因此，完整证明必须分别说明：最优解为何存在、为何不是角点、为何一阶条件足以刻画全局最优。仅写出
\[
\frac{\alpha}{1-\alpha}\frac{x_2}{x_1}=\frac{p_1}{p_2}
\]
只能得到候选比例，不能独立构成需求定理。

\begin{proposition}[严格单调与财富耗尽的再核对]
若 $u$ 严格单调，预算集为非空紧集，且所有价格严格为正，则每个最优束均耗尽预算；若 $u$ 仅弱单调，结论可能失败。
\end{proposition}
\begin{proof}
严格单调情形可由增加一个商品的微小扰动证明。弱单调反例取 $u(x)=\min\{x_1,1\}$，当 $x_1\ge1$ 后增加商品一不再提高效用，因此存在未耗尽预算的最优束。这个反例说明“预算线上的最优”来自严格单调假设，而不是来自预算约束本身。
\end{proof}

\begin{example}[同一偏好下的不同数值表示]
若 $u(x)=x_1+x_2$，则 $v(x)=3u(x)+7$ 和 $w(x)=\log(1+e^{u(x)})$ 都保持相同的偏好排序。三者的水平集相同，但函数值差、边际效用和曲率不同。因此，看到某个表示的导数符号时，必须先问该结论是否在严格递增变换下保持；偏好顺序保持，而边际效用水平、效用差值和曲率通常不保持。
\end{example}


\end{document}
